% ======================================================================
%  LE SYSTÈME CONVOLUTIF — RAPPORTS TYPIQUE ET NON TYPIQUES 1/k
%  Géométrie du spectre des nombres premiers — Théorie « L'Univers est au carré »
%  Auteur : Philippe Thomas Savard
% ======================================================================
\documentclass[11pt,a4paper]{article}

\usepackage[utf8]{inputenc}
\usepackage[T1]{fontenc}
\usepackage[french]{babel}
\usepackage{lmodern}
\usepackage{microtype}
\usepackage[a4paper,margin=2.4cm,headheight=15pt]{geometry}
\usepackage{amsmath,amssymb,amsfonts}
\usepackage[table]{xcolor}
\usepackage{booktabs,longtable,array,multirow}
\usepackage{fancyhdr}
\usepackage{lastpage}
\usepackage{titlesec}
\usepackage{enumitem}
\usepackage{fancyvrb}
\usepackage{tikz}
\usepackage{graphicx}
\usetikzlibrary{arrows.meta,positioning,shapes.geometric,calc,backgrounds}
\usepackage[most]{tcolorbox}
\usepackage[hidelinks]{hyperref}

\allowdisplaybreaks[3]
\setlength{\emergencystretch}{2em}
\setlength{\tabcolsep}{4pt}

% ---------- Palette « géométrie du spectre » ----------
\definecolor{spectralbleu}{RGB}{18,42,90}
\definecolor{spectralazur}{RGB}{41,98,168}
\definecolor{spectralor}{RGB}{196,148,34}
\definecolor{spectralviolet}{RGB}{96,54,140}
\definecolor{spectralgris}{RGB}{244,246,250}
\definecolor{spectralrouge}{RGB}{158,42,52}

% ---------- En-tête / pied de page ----------
\pagestyle{fancy}
\fancyhf{}
\fancyhead[L]{\small\color{spectralbleu}\textbf{www.universestaucarre.com} --- Dépôt GitHub de l'agent Gabriel multiloop local}
\fancyhead[R]{\small\color{spectralbleu}Page \thepage\ sur \pageref{LastPage}}
\renewcommand{\headrulewidth}{0.8pt}
\renewcommand{\headrule}{\color{spectralor}\hrule width\headwidth height 0.8pt}
\fancyfoot[C]{\small\color{spectralbleu}\itshape La géométrie du spectre des nombres premiers --- Le système convolutif}

% ---------- Titres ----------
\titleformat{\section}{\Large\bfseries\color{spectralbleu}}{\thesection.}{0.6em}{}[{\color{spectralor}\titlerule[1.2pt]}]
\titleformat{\subsection}{\large\bfseries\color{spectralazur}}{\thesubsection.}{0.6em}{}
\titleformat{\subsubsection}{\normalsize\bfseries\color{spectralviolet}}{\thesubsubsection.}{0.6em}{}

% ---------- Boîtes ----------
\newtcolorbox{boiteformule}[1][]{colback=spectralgris,colframe=spectralbleu,
  title={#1},fonttitle=\bfseries\small,fontupper=\small,
  enhanced,breakable,left=2mm,right=2mm,arc=2mm,boxrule=0.8pt}
\newtcolorbox{boitenote}[1][Note]{colback=white,colframe=spectralor,
  title={#1},fonttitle=\bfseries\small,fontupper=\small,
  enhanced,breakable,left=2mm,right=2mm,arc=2mm,boxrule=0.8pt}
\newtcolorbox{boitealerte}[1][Point de vigilance]{colback=white,colframe=spectralrouge,
  title={#1},fonttitle=\bfseries\small,fontupper=\small,
  enhanced,breakable,left=2mm,right=2mm,arc=2mm,boxrule=0.8pt}

\newcommand{\SA}{\mathrm{S}_A}
\newcommand{\SB}{\mathrm{S}_B}

\begin{document}
% ======================================================================
%  PAGE DE TITRE — style autodidacte mathématique, thème spectral
% ======================================================================
\begin{titlepage}
\begin{tikzpicture}[remember picture,overlay]
  \fill[spectralbleu] (current page.north west) rectangle ([yshift=-3.2cm]current page.north east);
  \fill[spectralor]   ([yshift=-3.2cm]current page.north west) rectangle ([yshift=-3.35cm]current page.north east);
  \fill[spectralbleu] (current page.south west) rectangle ([yshift=1.6cm]current page.south east);
  \fill[spectralor]   ([yshift=1.6cm]current page.south west) rectangle ([yshift=1.75cm]current page.south east);
  % Spectre des nombres premiers : graduations 1..61, barres dorées aux premiers
  \begin{scope}[shift={([yshift=-10.9cm]current page.north west)},x=2.5mm]
    \draw[spectralbleu!40,thick] (2,0) -- (61,0);
    \foreach \n in {2,3,...,61}{ \draw[spectralbleu!35] (\n,0) -- (\n,0.12); }
    \foreach \p in {2,3,5,7,11,13,17,19,23,29,31,37,41,43,47,53,59,61}{
      \draw[spectralor,line width=1.4pt] (\p,0) -- (\p,0.85);
      \fill[spectralor] (\p,0.85) circle (0.9pt); }
    \draw[spectralviolet,dashed,line width=1pt] (31.5,-0.25) -- (31.5,1.5)
      node[above]{\footnotesize\color{spectralviolet}$R_{sP}=\tfrac{1}{2}$};
  \end{scope}
  % Carré « L'Univers est au carré » avec hypoténuse
  \begin{scope}[shift={([xshift=15.2cm,yshift=-14.2cm]current page.north west)}]
    \draw[spectralazur,thick] (0,0) rectangle (2.4,2.4);
    \draw[spectralor,thick] (0,0) -- (2.4,2.4) node[midway,above left,sloped]{\footnotesize $\sqrt{a^{2}+b^{2}}$};
    \node[below] at (1.2,0) {\footnotesize $b=k^{i}$};
    \node[left] at (0,1.2) {\footnotesize $a=k^{i-1}$};
  \end{scope}
  \node[anchor=west,text=white,font=\small\bfseries] at ([xshift=2.2cm,yshift=-1.1cm]current page.north west)
    {www.universestaucarre.com};
  \node[anchor=west,text=white!85!spectralbleu,font=\footnotesize] at ([xshift=2.2cm,yshift=-1.75cm]current page.north west)
    {Dépôt GitHub de l'agent Gabriel multiloop local --- \texttt{agent-multiloop-Gabriel-local}};
  \node[anchor=west,text=white!85!spectralbleu,font=\footnotesize] at ([xshift=2.2cm,yshift=-2.35cm]current page.north west)
    {Théorie unifiée : \emph{L'Univers est au carré} --- Chapitre : \emph{La géométrie du spectre des nombres premiers}};
  \node[anchor=west,text=spectralbleu,font=\Huge\bfseries,align=left] at ([xshift=2.0cm,yshift=-6.6cm]current page.north west)
    {Le système convolutif};
  \node[anchor=west,text=spectralazur,font=\LARGE\bfseries] at ([xshift=2.0cm,yshift=-7.9cm]current page.north west)
    {Rapports typique et non typiques : $1/k = 1/2$ et $1/k \neq 1/2$};
  \node[anchor=west,text=spectralviolet,font=\large\itshape] at ([xshift=2.0cm,yshift=-8.9cm]current page.north west)
    {Système équationnel convolutif pour les requêtes multi-rapports du pipeline cognitif};
  \node[anchor=north west,align=left,font=\small] at ([xshift=2.0cm,yshift=-17.6cm]current page.north west)
    {\textbf{\color{spectralbleu}Auteur :} Philippe Thomas Savard\\[2pt]
     \textbf{\color{spectralbleu}Collaborateurs :} Copilot Microsoft --- Gordon Docker Desktop --- E1 emergent.sh\\
     \hspace{3.1cm} API Anthropic --- API OpenAI --- Cline --- Gemini Google\\[6pt]
     \textbf{\color{spectralbleu}Adresse :} 4-5354 Rue du Menuet, Lévis, Chaudière-Appalaches, Canada, G6X 2Y6\\
     \textbf{\color{spectralbleu}Courriel :} philippethomassavard@gmail.com\\[6pt]
     \textbf{\color{spectralbleu}Lieu et date :} Lévis, Chaudière-Appalaches, Canada --- le vingt-six septembre deux mille vingt-six};
  \node[anchor=south,text=white,font=\footnotesize] at ([yshift=0.6cm]current page.south)
    {Licence Apache 2.0 --- Document fourni tel quel, sans garantie --- Créé sans financement d'aucune nature};
\end{tikzpicture}
\phantom{x}
\end{titlepage}

% ======================================================================
%  AVERTISSEMENT ET LICENCE (texte du document original, orthographe corrigée)
% ======================================================================
\section*{Avertissement et licence}
\addcontentsline{toc}{section}{Avertissement et licence}

Les informations contenues dans le document que vous consultez sont sous
\textsc{Licence Apache~2.0}. Le document est présent sous la même licence sur
GitHub dans le dépôt public :
\url{https://github.com/2racinede4carreunivers-dev/agent-multiloop-Gabriel-local.git}.
Il est possible, selon certaines restrictions en lien avec la \textsc{Licence
Apache~2.0} et la réglementation du dépôt, de contribuer via PR, de cloner,
partager et modifier les informations présentes dans ce document. Le document
est fourni tel quel, sans garantie ni assurance quant à sa validité, et son
utilisation est permise dans la mesure de cet avertissement en ce qui concerne
sa validité. Le document a été créé sans financement de toute nature, ni
avantage, ni intérêt.

Le présent document est une section de la théorie unifiée \og L'Univers est au
carré \fg{}, et fait partie du chapitre de cette théorie unifiée : \og La
géométrie du spectre des nombres premiers \fg{}. Cette géométrie, incluse dans
la théorie mentionnée, est un outil dynamique conçu entièrement par son auteur
Philippe Thomas Savard en collaboration avec les entités mentionnées. Cet outil
dynamique est conçu pour résoudre la question de l'énigme de Bernhard Riemann
et la conjecture de la fonction Zêta de Riemann.

Le dépôt mathlib : \url{https://github.com/leanprover-community/mathlib4.git}
est utilisé pour certaines informations selon les termes de la licence du dépôt
en question. Les informations utilisées, pouvant avoir été modifiées, sont
incluses dans le fichier \texttt{.thy} de contre-validation du dépôt :
\url{agent-multiloop-Gabriel-local/theories/validation_zeta_lean.thy} ; la
licence du dépôt mathlib est incluse dans le fichier \texttt{.thy}. Cette
contre-validation, et la contre-validation jumelle dans le même sens
\texttt{validation\_hol\_unifiee.thy}, concernent toutes deux la validation
principale \texttt{methode\_spectral.thy}, présentes aux adresses du dépôt :
\url{agent-multiloop-Gabriel-local/theories/validation_hol_unifiee.thy} et
\url{agent-multiloop-Gabriel-local/theories/methode_spectral.thy},
disponibles en sept langues différentes sur le dépôt ; quelques annotations
sont en langue phonétique.

\tableofcontents
\newpage

% ======================================================================
\section{Rapports typique et non typiques : $1/k = 1/2$ et $1/k \neq 1/2$}
% ======================================================================

\subsection{Exemple de suites géométriques A et B}

Pour le rapport typique $t = 1/k = 1/2$, les suites A et B sont construites
selon le modèle de l'hypoténuse $\sqrt{a^{2}+b^{2}}$, dix termes par suite
(convention de notation : $(t^{m})^{2}$ signifie le carré de $t^{m}$ ; ainsi
$(t^{3})^{2} = t^{6}$) :

\begin{boiteformule}[Somme suite A --- modèle hypoténuse]
\[
\begin{aligned}
\SA ={}& \bigl((1^{1})^{2} + (t^{1})^{2}\bigr)^{1/2}
 + \bigl((t^{1})^{2} + (t^{2})^{2}\bigr)^{1/2}
 + \bigl((t^{2})^{2} + (t^{3})^{2}\bigr)^{1/2}
 + \bigl((t^{3})^{2} + (t^{4})^{2}\bigr)^{1/2} \\
& + \bigl((t^{4})^{2} + (t^{5})^{2}\bigr)^{1/2}
 + \bigl((t^{5})^{2} + (t^{6})^{2}\bigr)^{1/2}
 + \bigl((t^{6})^{2} + (t^{7})^{2}\bigr)^{1/2}
 + \bigl((t^{7})^{2} + (t^{8})^{2}\bigr)^{1/2} \\
& + \bigl((t^{8} - t^{6})^{2} + (t^{9} - t^{7})^{2}\bigr)^{1/2}
 + \bigl((t^{9} - t^{7})^{2} + (t^{10} - t^{8})^{2}\bigr)^{1/2}.
\end{aligned}
\]
\end{boiteformule}

\begin{boiteformule}[Somme suite B --- modèle hypoténuse (saut structurel en position 6)]
\[
\begin{aligned}
\SB ={}& \bigl((1^{1})^{2} + (t^{1})^{2}\bigr)^{1/2}
 + \bigl((t^{1})^{2} + (t^{2})^{2}\bigr)^{1/2}
 + \bigl((t^{2})^{2} + (t^{3})^{2}\bigr)^{1/2} \\
& + \bigl((t^{3})^{2} + (t^{4})^{2}\bigr)^{1/2}
 + \bigl((t^{4})^{2} + (t^{5})^{2}\bigr)^{1/2} \\
& + \bigl((t^{6})^{2} + (t^{7})^{2}\bigr)^{1/2}
 + \bigl((t^{7})^{2} + (t^{8})^{2}\bigr)^{1/2}
 + \bigl((t^{8})^{2} + (t^{9})^{2}\bigr)^{1/2} \\
& + \bigl((t^{10} - t^{8})^{2} + (t^{11} - t^{9})^{2}\bigr)^{1/2} \\
& + \bigl((t^{11} - t^{9})^{2} + (t^{12} - t^{10})^{2}\bigr)^{1/2}.
\end{aligned}
\]
\end{boiteformule}

\begin{boitenote}[Décalage structurel A $\to$ B]
La suite B ne comporte pas de terme $\bigl((t^{5})^{2}+(t^{6})^{2}\bigr)^{1/2}$ :
c'est le \emph{saut Zêta}, décalage structurel de A vers B. Les deux suites sont
identiques pour les positions $i \le 5$, puis divergent à partir de $i = 6$.
\end{boitenote}

\subsection{Reconstruction du nombre premier pour $n = 10$}

\subsubsection{Quatre possibilités pour le Digamma}

\begin{enumerate}[nosep]
  \item Additionner la 8\textsuperscript{ième} position de la suite A à la somme de la suite A 10 termes ($n=10$).
  \item Soustraire la 8\textsuperscript{ième} position de la suite A à la somme suite A 10 termes ($n=10$).
  \item Additionner la 7\textsuperscript{ième} position de la suite A à la somme de la suite A 10 termes ($n=10$).
  \item Soustraire la 7\textsuperscript{ième} position de la suite A à la somme de la suite A 10 termes ($n=10$).
\end{enumerate}

Les quatre possibilités permettent de déterminer la valeur du Digamma calculé.

\subsection{Généralisation des équations des suites A et B}

\textbf{1. Coefficient A :}
\[
\frac{\text{Somme suite A}(10\ \text{termes}) - \text{Somme suite A}(9\ \text{termes})}{t^{8}}
= \text{Coefficient A}.
\]

\textbf{Coefficient B :}
\[
\frac{\text{Somme suite B}(10\ \text{termes}) - \text{Somme suite B}(9\ \text{termes})}{t^{8}}
= \text{Coefficient B}.
\]

\textbf{2. Reconstruire les suites A et B en déterminant la somme des suites A
et B pour toutes les valeurs $n$ :}
\[
\frac{\text{Coefficient A}}{x} \times k^{n} - \text{Reste} = \text{Somme suite A},
\qquad
\frac{\text{Coefficient A}}{\text{Somme suite A}}\, k^{10} = \text{Reste} + x .
\]
\[
(\text{Reste} + x) - \text{la valeur entière} = x .
\]

La valeur de $(\text{Reste} + x)$ forme deux blocs A et B : A étant la partie
entière de la valeur $(\text{Reste} + x)$, et B formant la partie $< 1$ qui est
décimale de $(\text{Reste} + x)$.
\begin{align*}
\frac{\text{Coefficient A}}{x}\times k^{10}-\text{Reste}
&=\text{Somme suite A}+\text{Reste},\\
\text{Somme suite A}-(\text{Somme suite}+\text{Reste})
&=\text{Reste}.
\end{align*}

\begin{boiteformule}[Équation reconstruite A généralisant la valeur de $n$]
\[
\frac{\text{Coefficient A}}{x} \times k^{n} - \text{Reste} = \text{Somme suite A},
\qquad \text{quand } n \text{ est un entier strictement positif.}
\]
\end{boiteformule}

De même pour la suite B :
\[
\frac{\text{Coefficient B}}{x} \times k^{n} - \text{Reste} = \text{Somme suite B},
\qquad
\frac{\text{Coefficient B}}{\text{Somme suite B}}\, k^{10} = \text{Reste} + x ,
\]
\begin{align*}
\frac{\text{Coefficient B}}{x}\times k^{10}-\text{Reste}
&=\text{Somme suite B}+\text{Reste},\\
\text{Somme suite B}-(\text{Somme suite}+\text{Reste})
&=\text{Reste}.
\end{align*}

\begin{boiteformule}[Équation reconstruite B généralisant la valeur de $n$]
\[
\frac{\text{Coefficient B}}{x} \times k^{n} - \text{Reste} = \text{Somme suite B},
\qquad \text{quand } n \text{ est un entier strictement positif.}
\]
\end{boiteformule}

\subsection{Approche généralisée de niveau 1}

\begin{boiteformule}[Sommes des suites A et B (forme formelle, $n = 10$)]
\begin{align*}
t^{1}+t^{2}+t^{3}+t^{4}+t^{5}+t^{6}+t^{7}+t^{8}+(t^{9}-t^{7})+(t^{10}-t^{8}) &= \text{Somme suite A},\\
t^{1}+t^{2}+t^{3}+t^{4}+t^{5}+t^{7}+t^{8}+t^{9}+(t^{10}-t^{8})+(t^{11}-t^{9}) &= \text{Somme suite B}.
\end{align*}
\end{boiteformule}

Digamma, 4 possibilités :
\begin{enumerate}[nosep]
  \item $-$ 8\textsuperscript{ième} position suite A : $-t^{8}$
  \item $+$ 8\textsuperscript{ième} position suite A : $+t^{8}$
  \item $-$ 7\textsuperscript{ième} position suite A : $-t^{7}$
  \item $+$ 7\textsuperscript{ième} position suite A : $+t^{7}$
\end{enumerate}

Digamma calculé : $\ \text{Somme suite A} - t^{8}$ ;\quad
$\text{Somme suite A} + t^{8}$ ;\quad
$\text{Somme suite A} - t^{7}$ ;\quad
$\text{Somme suite A} + t^{7}$.

\subsubsection{Détermination du nombre premier}
\begin{boiteformule}[Reconstruction du nombre premier]
\[
\frac{\text{Somme suite B} - \text{Digamma calculé}}{6\textsuperscript{ième}\ \text{position suite A}\ (\text{Zêta})}
= \text{Nombre premier}.
\]
\end{boiteformule}

\subsubsection{Équations des sommes des suites A et B}

Déterminer le coefficient A :
$\displaystyle \frac{(\text{Somme suite A}_{n=10}) - \text{Somme suite A}_{n=9})}{t^{8}}$ \quad
Déterminer le coefficient B :
$\displaystyle \frac{(\text{Somme suite B}_{n=10}) - \text{Somme suite B}_{n=9})}{t^{8}}$

Déterminer l'équation A :
\begin{align*}
\frac{\text{Coefficient A}}{\text{Somme suite A}_{n=10}}\,t^{10}
&=x+\text{Reste},\\
(x+\text{Reste})-\text{Valeur réelle}&=x,\\
\frac{\text{Coefficient A}}{x}\times t^{10}
&=\text{Somme suite A}+\text{Reste},\\
(\text{Somme suite A}+\text{Reste})-(\text{Somme suite A})
&=\text{Reste}.
\end{align*}

\textbf{Équation A :} $\displaystyle \frac{\text{Coefficient A}}{x} \times t^{n} - \text{Reste} = \text{Somme suite A}$,
quand $n$ est un entier strictement positif.

Déterminer l'équation B :
\begin{align*}
\frac{\text{Coefficient B}}{\text{Somme suite B}_{n=10}}\,t^{10}
&=x+\text{Reste},\\
(x+\text{Reste})-\text{Valeur réelle}&=x,\\
\frac{\text{Coefficient B}}{x}\times t^{10}
&=\text{Somme suite B}+\text{Reste},\\
(\text{Somme suite B}+\text{Reste})-(\text{Somme suite B})
&=\text{Reste}.
\end{align*}

\textbf{Équation B :} $\displaystyle \frac{\text{Coefficient B}}{x} \times t^{n} - \text{Reste} = \text{Somme suite B}$,
quand $n$ est un entier strictement positif.

\begin{boitenote}[Inspiration]
Cette équation est tirée de la formule déjà connue des suites géométriques :
\[
\text{Si } q \neq 1,\quad S = 1 + q + q^{2} + q^{3} \dots q^{n} = \frac{1 - q^{n+1}}{1 - q}.
\]
\end{boitenote}

% ======================================================================
\section{Définition et propriétés des rapports non typiques}
% ======================================================================

Les rapports non typiques $1/k$ sont définis pour toute valeur $1/k < 1/2$,
notamment $1/3$, $1/4$, $1/5$, $1/6$, \dots, $1/k$.

Contrairement au rapport typique $1/2$, la valeur de $n$ ne correspond pas à la
position du nombre premier. Elle représente exclusivement le nombre de termes
des suites A et B.

\subsection{Exemple pour $n=10$ des suites A et B pour un rapport non typique $1/k = 1/3$}

\begin{boiteformule}[Sommes des suites A et B --- rapport $1/3$, $n = 10$]
\begin{align*}
3^{1}+3^{2}+3^{3}+3^{4}+3^{5}+3^{6}+3^{7}+3^{8}+(3^{9}-3^{7})+(3^{10}-3^{8})
  &= \text{Somme suite A} = 79\,824,\\
3^{1}+3^{2}+3^{3}+3^{4}+3^{5}+3^{7}+3^{8}+3^{9}+(3^{10}-3^{8})+(3^{11}-3^{9})
  &= \text{Somme suite B} = 238\,746.
\end{align*}
\end{boiteformule}

$\text{Digamma} = 3^{8}$ (8\textsuperscript{ième} position, soustraite) ;\quad
Digamma calculé $= 79\,824 - 3^{8} = 73\,263$.

Déterminer le nombre premier :
\[
\frac{238\,746 - 73\,263}{3^{6}} = 227 \quad (49\textsuperscript{ième}\ \text{nombre premier}).
\]

\subsubsection{Rapport non typique $1/k = 1/3$ --- équations généralisées}

Somme suite A$(227) = 79\,824$.

\textbf{1.} Déterminer la somme de la suite B du nombre précédent 227, qui est 223 :
\begin{align*}
\text{Somme suite A}(227)-3^{6}
  &= \text{Somme suite B}(223),\\
\text{Somme suite B}(223)
  &= 79\,824-3^{6}=79\,095.
\end{align*}
9 termes, suite B(223) :
\begin{align*}
3^{1}+3^{2}+3^{3}+3^{4}+3^{5}+3^{7}+3^{8}
&+(3^{9}-3^{7})+(3^{10}+3^{8})\\
&=79\,095.
\end{align*}
9 termes, suite A(223) :
\begin{align*}
3^{1}+3^{2}+3^{3}+3^{4}+3^{5}+3^{6}+3^{7}
&+(3^{8}-3^{6})+(3^{9}+3^{7})\\
&=26\,607.
\end{align*}

\textbf{2.} Déterminer le coefficient B, avec Somme suite B$(227) = 238\,746$ :
\begin{align*}
\frac{\text{Somme suite B}(227) - \text{Somme suite B}(223)}{3^{8}}
&= \text{Coefficient B},\\
\frac{238\,746 - 79\,095}{3^{8}} &= \frac{219}{9} = \text{Coefficient B}.
\end{align*}
\begin{boitenote}
Il s'agit de l'écart le plus petit entre deux suites B pour toutes les valeurs
de $n$ positives.
\end{boitenote}

\textbf{3.} Déterminer l'écart minimum entre les termes $n=1$ et $n=2$ pour le
rapport non typique $1/k = 1/3$ des suites B :
\begin{align*}
\frac{\text{Somme suite B}(227) - \text{Somme suite B}(223)}{3^{6}}
&= \text{différence entre les deux suites B}(227)\ \text{et B}(223),\\
\frac{238\,746 - 79\,095}{3^{6}} &= 219.
\end{align*}

\textbf{4.} Écart entre les suites A(227) et B(223) :
\begin{align*}
\frac{219}{3} &= 73,\\
\left(\frac{79\,824}{3^{6}} - 73\right) \times 3^{6}
&= 26\,607 \quad \text{Somme suite A}(223).
\end{align*}

\textbf{5.} Déterminer le coefficient pour l'équation de la suite A :
\begin{align*}
\frac{\text{Somme suite A}(227) - \text{Somme suite A}(223)}{3^{8}}
&= \text{Coefficient A},\\
\frac{79\,824 - 26\,607}{3^{8}} &= \frac{73}{9} = \text{Coefficient A}.
\end{align*}

\textbf{6.} Équations pour les sommes des suites A et B --- l'équation est
inspirée de celle déjà connue concernant les suites géométriques
($S = \frac{1-q^{n+1}}{1-q}$, $q \neq 1$). Déterminer l'équation pour la
suite A, coefficient A $= 73/9$ :
\begin{align*}
\left(\frac{73}{9}\,x \times 3^{10}\right)-\text{Reste} &= 79\,824,\\
\frac{73}{9}\Big/\frac{238\,746+\text{reste}}{3^{10}}
&\colon \frac{73}{9}/1.351826449 = 6.000112748,\\
6.000112748-0.000112748 &= 6=x,\\
\frac{73}{9}\times6\times3^{10}-\text{Reste} &= 79\,824,\\
\frac{73}{9}\times6\times3^{10} &= 79\,825.5,\\
79\,825.5-79\,824 &= \frac{3}{2}=\text{Reste}.
\end{align*}

\begin{boiteformule}[Équation somme suite A --- rapport $1/3$]
\[
\text{Équation } n=10 :\quad \left(\frac{73}{9}\times 6 \times 3^{10}\right) - \frac{3}{2} = 79\,824.
\]
\[
\text{Forme générale :}\quad \left(\frac{73}{9}\times 6 \times 3^{n}\right) - \frac{3}{2} = \text{Somme suite A}
\quad (n \text{ entier strictement positif}),
\]
\[
\frac{73}{9}\times 3^{-n} - \frac{3}{2} = \text{Somme suite A négative}
\quad (n \text{ entier strictement négatif}).
\]
\end{boiteformule}

\textbf{7.} Déterminer l'équation pour la suite B (coefficient B $= 219/9$) :
\[
\left(\frac{219}{9}\, x \times 3^{10}\right) - \text{Reste} = 238\,746,
\qquad
\frac{219}{9}\times 4.043184474 = 6.018358423,
\]
\[
6.018358423 - 0.018358423 = 6 = x,
\qquad
\frac{219}{9}\times 6 \times 3^{10} - \text{Reste} = 238\,746,
\]
\[
\frac{219}{9}\times 6 \times 3^{10} = 238\,476.5,
\qquad
238\,476.5 - 238\,746 = 730.5 = \text{Reste},
\qquad
\frac{730.5}{3/2} = 487.
\]

\begin{boiteformule}[Équation somme suite B --- rapport $1/3$]
\[
\text{Équation } n=10 :\quad \left(\frac{219}{9}\times 6 \times 3^{10}\right) - \frac{3}{2}(487) = 238\,746.
\]
\[
\text{Forme générale :}\quad \left(\frac{219}{9}\times 6 \times 3^{n}\right) - \frac{3}{2}(487) = \text{Somme suite B}
\quad (n \text{ entier strictement positif}),
\]
\[
\frac{219}{9}\times 3^{-n} - \frac{3}{2}(487) = \text{Somme suite B négative}
\quad (n \text{ entier strictement négatif}).
\]
\end{boiteformule}

\textbf{9.} Caractéristiques intrinsèques de la valeur de $n$ pour les rapports
non typiques $1/k \neq 1/2$ : la valeur de $n$ pour des suites A et B non
typiques ne correspond pas à la position du nombre premier ; la valeur de $n$
correspond à la quantité de termes dans les suites A et B, tant pour les suites
positives que négatives.

\subsection{Rapport non typique $1/k = 1/4$ ($n = 10$ ; $p = 947$, 161\textsuperscript{ème} nombre premier)}

\begin{boiteformule}[Sommes des suites A et B --- rapport $1/4$, $n = 10$]
\begin{align*}
4^{1}+4^{2}+4^{3}+4^{4}+4^{5}+4^{6}+4^{7}+4^{8}+(4^{9}-4^{7})+(4^{10}-4^{8})
  &= \text{Somme suite A} = 1\,316\,180,\\
4^{1}+4^{2}+4^{3}+4^{4}+4^{5}+4^{7}+4^{8}+4^{9}+(4^{10}-4^{8})+(4^{11}-4^{9})
  &= \text{Somme suite B} = 5\,260\,628.
\end{align*}
\end{boiteformule}

Reconstruire le nombre premier --- Digamma : $+4^{8}$ :
\begin{align*}
\text{Digamma calculé : }1\,316\,180+4^{8}
&=1\,316\,180+65\,536=1\,381\,716,\\
\frac{5\,260\,628-1\,381\,716}{4^{6}}
&=\frac{3\,878\,912}{4096}=947.
\end{align*}

Suites A et B pour $n = 9$ termes (941) :
\begin{align*}
4^{1}+4^{2}+4^{3}+4^{4}+4^{5}+4^{6}+4^{7}+(4^{8}-4^{6})+(4^{9}-4^{7}) &= 329\,044,\\
4^{1}+4^{2}+4^{3}+4^{4}+4^{5}+4^{7}+4^{8}+(4^{9}-4^{7})+(4^{10}-4^{8}) &= 1\,312\,084.
\end{align*}

Déterminer Somme suite B(941) et les écarts :
\[
1316180 - 4096 = 1\,312\,084,
\qquad
\frac{5\,260\,628 - 1\,312\,084}{4^{6}} = \frac{3\,948\,544}{4096} = 964,
\qquad
\frac{964}{4} = 241,
\]
\begin{align*}
\frac{1\,316\,180 - 329\,044}{4^{6}}
&= \frac{987\,136}{4096} = 241,\\
\text{Coefficient A} &= \frac{987\,136}{65\,536} = \frac{241}{16},\\
\text{Coefficient B} &= \frac{3\,948\,544}{65\,536}
= \frac{964}{16} = \frac{241}{4}.
\end{align*}

\begin{boiteformule}[Équations généralisées --- rapport $1/4$]
\begin{align*}
\text{Équation A }(n=10):\quad
\left(\frac{241}{16}\times12\times4^{10}\right)-\frac{4}{3}
&=1\,316\,180,\\
\frac{241}{16}\Big/\frac{1\,316\,180}{4^{10}}
&=12.000001216\Rightarrow x=12,\\
\text{Équation A généralisée :}\quad
\left(\frac{241}{192}\times4^{n}\right)-\frac{4}{3}
&=\text{Somme suite A}\quad(n\in\mathbb{Z},\ n>0),\\
\frac{241}{16}\times4^{-n}-\frac{4}{3}
&=\text{Somme suite A négative}.
\end{align*}
\begin{align*}
\text{Équation B }(n=10):\quad
\left(\frac{241}{4}\times12\times4^{10}\right)-3073\left(\frac{4}{3}\right)
&=5\,260\,628,\\
\text{Équation B généralisée :}\quad
\left(\frac{241}{4}\times12\times4^{n}\right)-3073\left(\frac{4}{3}\right)
&=\text{Somme suite B},\\
\frac{241}{4}\times4^{n}-3073\left(\frac{4}{3}\right)
&=\text{Somme suite B négative}.
\end{align*}
\end{boiteformule}

\subsection{Rapport non typique $1/k = 1/5$}

\begin{boiteformule}[Sommes des suites A et B --- rapport $1/5$, $n = 10$]
\begin{align*}
5^{1}+5^{2}+5^{3}+5^{4}+5^{5}+5^{6}+5^{7}+5^{8}+(5^{9}-5^{7})+(5^{10}-5^{8})
  &= 11\,738\,280 \quad \text{Somme suite A},\\
5^{1}+5^{2}+5^{3}+5^{4}+5^{5}+5^{7}+5^{8}+5^{9}+(5^{10}-5^{8})+(5^{11}+5^{9})
  &= 58\,675\,780 \quad \text{Somme suite B}.
\end{align*}
\end{boiteformule}

\textbf{1.} Reconstruire le nombre premier --- Digamma : $+$7\textsuperscript{ième}
position suite A : $+5^{7}$ :
\[
\text{Digamma calculé : } 11\,738\,280 + 5^{7} = 11\,816\,405,
\qquad
\frac{58\,675\,780 - 11\,816\,405}{5^{6}} = 2999.
\]

\textbf{2.} Suites A et B pour $n = 9$ termes (2993) :
\begin{align*}
5^{1}+5^{2}+5^{3}+5^{4}+5^{5}+5^{6}+5^{7}+(5^{8}-5^{6})+(5^{9}-5^{7}) &= 2\,347\,655,\\
5^{1}+5^{2}+5^{3}+5^{4}+5^{5}+5^{7}+5^{8}+(5^{9}-5^{7})+(5^{10}-5^{8}) &= 11\,722\,655.
\end{align*}

\textbf{3.} Déterminer les équations servant de manière généralisée à
l'ensemble des valeurs de $n$ entières positives et négatives pour les suites A
et B.

\textbf{4.} Déterminer Somme suite B(2993) :
$\ \text{Somme suite B}(2999) - 5^{6} = \text{Somme suite A}$ ;
$\ 11\,738\,280 - 5^{6} = 11\,722\,655$.

\textbf{5.} $\displaystyle \frac{\text{Somme suite B}(2999) - \text{Somme suite B}(2993)}{5^{6}}
= \frac{58\,675\,780 - 11\,722\,655}{5^{6}} = 3005$ \quad (écart entre les suites B $= 3005$).

\textbf{6.} $\displaystyle \frac{3005}{5} = 601$ (écart entre les suites A(2999) et A(2993)) ;
$\displaystyle \frac{11\,738\,280 - 2\,347\,655}{5^{6}} = 601$.

\textbf{7.} Coefficients : \quad
Coefficient A $= \dfrac{11\,738\,280 - 2\,347\,655}{5^{8}} = \dfrac{601}{25}$ ;\qquad
Coefficient B $= \dfrac{58\,675\,780 - 11\,722\,655}{5^{8}} = \dfrac{601}{5}$.

\textbf{8.} Équation A :
\[
\left(\frac{601}{25}\, x \times 5^{10}\right) - \text{Reste} = 11\,738\,280,
\qquad
\frac{601}{25}\Big/\frac{11\,738\,280}{5^{10}} = 20.00000213 \Rightarrow x = 20,
\]
\[
\frac{601}{25}\times 20 \times 5^{10} - \text{Reste} = 11\,738\,281.25,
\qquad
11\,738\,280 - 11\,738\,281.25 = -\frac{5}{4}.
\]

\begin{boiteformule}[Équations généralisées --- rapport $1/5$]
\begin{align*}
\text{Équation A }(n=10):\quad
\left(\frac{601}{25}\times20\times5^{10}\right)-\frac{5}{4}
&=11\,738\,280,\\
\left(\frac{601}{25}\times20\times5^{n}\right)-\frac{5}{4}
&=\text{Somme suite A},\\
\text{Équation A négative :}\quad
\frac{601}{25}\times5^{n}-\frac{5}{4}
&=\text{Somme suite A},\\
\text{Équation B :}\quad
\frac{601}{5}\times20\times5^{10}-\text{Reste}
&=58\,675\,780,\\
\frac{601}{5}\Big/\frac{58\,675\,780}{5^{10}}
&=20.13847323\Rightarrow x=20,\\
\left(\frac{601}{5}\times20\times5^{10}\right)-\text{Reste}
&=58\,691\,406.25,\\
58\,691\,406.25-58\,675\,780 &=15\,626.25,\\
\frac{15\,626.25}{1.25} &=12\,501,\\
\text{Équation B }(n=10):\quad
\frac{601}{5}\times20\times5^{10}
-12\,501\left(\frac{5}{4}\right)
&=58\,675\,780,\\
\frac{601}{5}\times20\times5^{n}
-12\,501\left(\frac{5}{4}\right)
&=\text{Somme suite B}.
\end{align*}
\end{boiteformule}

\subsection{Rapport non typique $1/k = 1/6$}

\begin{boiteformule}[Sommes des suites A et B --- rapport $1/6$, $n = 10$]
\begin{align*}
6^{1}+6^{2}+6^{3}+6^{4}+6^{5}+6^{6}+6^{7}+6^{8}+(6^{9}-6^{7})+(6^{10}+6^{8})
  &= \text{Somme suite A} = 70\,599\,858,\\
6^{1}+6^{2}+6^{3}+6^{4}+6^{5}+6^{7}+6^{8}+6^{9}+(6^{10}-6^{8})+(6^{11}+6^{9})
  &= \text{Somme suite B} = 423\,552\,498.
\end{align*}
\end{boiteformule}

Digamma $= +6^{8}$ --- la septième position est la position pour déterminer le
Digamma calculé et est soustraite :
\begin{align*}
\text{Digamma calculé} &=70\,599\,858+6^{7}=72\,279\,474,\\
\frac{423\,552\,498-72\,279\,474}{6^{6}}
&=7529\quad\text{nombre premier}.
\end{align*}

\textbf{1.} Suites A et B de 9 termes (7523) :
\begin{align*}
\begin{gathered}
6^{1}+6^{2}+6^{3}+6^{4}+6^{5}\\
{}+6^{6}+6^{7}+(6^{8}-6^{6})+(6^{9}-6^{7})
\end{gathered}
&=11\,766\,642\quad\text{Somme suite A 9 termes (7523)},\\
\begin{gathered}
6^{1}+6^{2}+6^{3}+6^{4}+6^{5}\\
{}+6^{7}+6^{8}+(6^{9}-6^{7})+(6^{10}-6^{8})
\end{gathered}
&=70\,553\,202\quad\text{Somme suite B 9 termes (7523)}.
\end{align*}
\[
\text{Somme suite A}(7529) - 6^{6} = \text{Somme suite B},
\qquad
70\,599\,858 - 6^{6} = 70\,553\,202.
\]

\textbf{2.} Écart des sommes des suites B :
$\displaystyle \frac{423\,552\,498 - 70\,553\,202}{6^{6}} = 7566$.

\textbf{3.} Écart entre les sommes des suites A(7529) et A(7523) :
$\displaystyle \frac{7566}{6} = 1261$.

\textbf{4.} Coefficient A :
$\displaystyle \frac{70\,599\,858 - 11\,766\,642}{6^{8}} = \frac{1261}{36}$.

\textbf{5.} Coefficient B :
$\displaystyle \frac{423\,552\,498 - 70\,553\,202}{6^{8}} = \frac{1261}{6}$.

\textbf{6.} Équations généralisant la valeur de $n$ pour les suites A et B du
rapport non typique $1/k = 1/6$ :
\begin{align*}
\frac{1261}{36}\, x \times 6^{10} - \text{Reste} = 70\,599\,858,
\qquad
\frac{1261}{36}\Big/\frac{70\,599\,858}{6^{10}}
&=30.00000051\Rightarrow x=30,\\
\frac{1261}{36}\times30\times6^{10}-\text{Reste}
&=70\,599\,859.2,\\
\frac{-46\,657.2}{6/5} &=-38\,881\times\frac{6}{5},\\
\frac{1261}{36}\times30\times6^{10}
-38\,881\left(\frac{6}{5}\right)
&=70\,553\,202.
\end{align*}
Pour la somme de la suite B :
\begin{align*}
\frac{1261}{6}\, x \times 6^{10} - \text{Reste} = 423\,552\,498,
\qquad
\frac{1261}{6}\Big/\frac{423\,552\,498}{6^{10}}
&=30.0033047\Rightarrow x=30,\\
\frac{1261}{6}\times30\times6^{10}-\text{Reste}
&=423\,599\,155.2,\\
\frac{-46\,657.2}{6/5} &=-38\,881,\\
\frac{1261}{6}\times30\times6^{10}-\text{Reste}
&=423\,552\,498.
\end{align*}

\subsection{Rapport non typique $1/k = 1/7$}

\begin{boiteformule}[Sommes des suites A et B --- rapport $1/7$, $n = 10$]
\begin{align*}
7^{1}+7^{2}+7^{3}+7^{4}+7^{5}+7^{6}+7^{7}+7^{8}+(7^{9}-7^{7})+(7^{10}-7^{8})
  &= \text{Somme suite A} = 322\,966\,112,\\
7^{1}+7^{2}+7^{3}+7^{4}+7^{5}+7^{7}+7^{8}+7^{9}+(7^{10}-7^{8})+(7^{11}-7^{9})
  &= \text{Somme suite B} = 2\,260\,645\,142.
\end{align*}
\end{boiteformule}

Reconstruire le nombre premier --- Digamma : $-7^{8}$ :
\begin{align*}
\text{Digamma calculé : }322\,966\,112-7^{8}
&=322\,966\,112-5\,764\,801=317\,201\,321,\\
\frac{2\,260\,645\,142-317\,201\,321}{7^{6}}
&=\frac{1\,943\,443\,821}{117\,649}=16\,519.
\end{align*}

Suites A et B pour $n = 9$ termes (16\,511) :
\begin{align*}
7^{1}+7^{2}+7^{3}+7^{4}+7^{5}+7^{6}+7^{7}+(7^{8}-7^{6})+(7^{9}-7^{7}) &= 46\,138\,018,\\
7^{1}+7^{2}+7^{3}+7^{4}+7^{5}+7^{7}+7^{8}+(7^{9}-7^{7})+(7^{10}-7^{8}) &= 322\,848\,463.
\end{align*}

Équations pour les suites A et B du rapport non typique $1/7$ :
\begin{align*}
\text{Somme suite A}(16519)-7^{6}
&=322\,966\,112-117\,649\\
&=322\,848\,463=\text{Somme suite B (9 termes)},\\
\frac{2\,260\,645\,142-322\,848\,463}{7^{6}}
&=\frac{1\,937\,796\,679}{117\,649}=16\,471,\\
\frac{16\,471}{7} &=2353,\\
\frac{322\,966\,112-46\,138\,018}{7^{6}}
&=\frac{276\,828\,094}{117\,649}=2353,\\
\text{Coefficient A}
&=\frac{276\,828\,094}{5\,764\,801}=\frac{2353}{49},\\
\text{Coefficient B}
&=\frac{1\,937\,796\,679}{5\,764\,801}=\frac{2353}{7}.
\end{align*}

\begin{boiteformule}[Équations généralisées --- rapport $1/7$]
\begin{align*}
\text{Équation A }(n=10):\quad
\left(\frac{2353}{49}\times42\times7^{10}\right)-\frac{7}{6}
&=322\,966\,112,\\
\frac{2353}{49}\Big/\frac{322\,966\,112}{7^{10}}
&=42.00000006\Rightarrow x=42,\\
\text{Forme générale A :}\quad
\left(\frac{2353}{2058}\times7^{n}\right)-\frac{7}{6}
&=\text{Somme suite A},\\
\frac{2353}{49}\times7^{-n}-\frac{7}{6}
&=\text{Somme suite A négative},\\
\text{Équation B }(n=10):\quad
\left(\frac{2353}{7}\times42\times7^{10}\right)
-100\,843\left(\frac{7}{6}\right)
&=2\,260\,645\,142,\\
\frac{705\,901}{6}\Big/\frac{7}{6} &=100\,843,\\
\text{Forme générale B :}\quad
\left(\frac{2353}{294}\times7^{n}\right)
-100\,843\left(\frac{7}{6}\right)
&=\text{Somme suite B}.
\end{align*}
\end{boiteformule}

% ======================================================================
\section{Système d'équations convolutif généralisé --- pour tout rapport $1/k$ ($k \ge 2$) et tout $n$ commun}
% ======================================================================

\subsection{Niveau 1 --- Constantes Savard universelles}

\begin{boiteformule}[Constantes universelles]
\[
\alpha_A(k) = \frac{2\,(k^{4} - k^{2} + 1)}{(k-1)\cdot k^{3}},
\qquad
\alpha_B(k) = k \cdot \alpha_A(k),
\]
\[
\delta_A(k) = \frac{k}{k-1},
\qquad
\delta_B(k) = \frac{k^{7} - k^{6} + k}{k-1}.
\]
\end{boiteformule}

\subsection{Niveau 2 --- Formules fermées ($n$ quelconque, $n$ commun à tous les rapports)}

\begin{boiteformule}[Formules fermées des sommes]
\[
\SA(k,n) = \frac{\alpha_A(k)}{2}\cdot k^{n} - \delta_A(k),
\qquad
\SB(k,n) = \frac{\alpha_B(k)}{2}\cdot k^{n} - \delta_B(k).
\]
\end{boiteformule}

\begin{boitenote}[Invariant spectral prouvé algébriquement (Théorème XIV.3)]
\[
\frac{\SA(k,n_1) - \SA(k,n_2)}{\SB(k,n_1) - \SB(k,n_2)} = \frac{1}{k}
\qquad \forall\, n_1 \neq n_2,\ k \ge 2.
\]
\end{boitenote}

\subsection{Niveau 3 --- Construction terme à terme (Suite A, $n \ge 8$)}

\begin{center}
\small
\begin{tabular}{@{}c l l@{}}
\toprule
\rowcolor{spectralgris}
Position $i$ & Terme $A(k,n,i)$ & Terme $B(k,n,i)$ \\
\midrule
$1 \le i \le n-2$ & $k^{i}$ & $k^{i}$ ($i \le 5$) ; $k^{7}$ ($i=6$, Saut Zêta) ; $k^{i+1}$ ($7 \le i \le n-2$) \\
$i = n-1$ & $k^{n-1} - k^{n-3}$ & $k^{n} - k^{n-2}$ \\
$i = n$ & $k^{n} - k^{n-2}$ & $k^{n+1} - k^{n-1}$ \\
\bottomrule
\end{tabular}
\end{center}

\subsection{Niveau 4 --- Mécanisme d'ancrage à $n = 10$ (4 candidats Digamma)}

Les quatre essais dans l'ordre déterministe :

\begin{center}
\small
\begin{tabular}{@{}c c c l l@{}}
\toprule
\rowcolor{spectralgris}
\# & Position & Signe & Digamma & Candidat $P$ \\
\midrule
1 & 8 & $-$ & $\SA(k,10) - k^{8}$ & $(\SB - \text{Digamma})/k^{6}$ \\
2 & 8 & $+$ & $\SA(k,10) + k^{8}$ & $(\SB - \text{Digamma})/k^{6}$ \\
3 & 7 & $+$ & $\SA(k,10) + k^{7}$ & $(\SB - \text{Digamma})/k^{6}$ \\
4 & 7 & $-$ & $\SA(k,10) - k^{7}$ & $(\SB - \text{Digamma})/k^{6}$ \\
\bottomrule
\end{tabular}
\end{center}

$\to$ Les quatre quotients sont examinés dans cet ordre. Si un seul candidat
est premier, il peut constituer l'ancre standard. Si plusieurs le sont, toutes
les branches sont conservées et aucune n'est choisie automatiquement ; une
branche inscrite au catalogue peut être utilisée seulement comme choix explicite
et doit rester distinguée d'une ancre unique. Si aucun candidat n'est premier,
aucune ancre standard n'est produite.
\[
P\bigl(1/k,\ n=10\bigr) = \frac{\SB(k,10) - \bigl[\SA(k,10) \pm k^{7}\ \text{ou}\ \pm k^{8}\bigr]}{k^{6}}.
\]

\subsection{Niveau 5 --- Reconstruction pour $n \neq 10$ ($n$ commun)}

\begin{boiteformule}[Loi de rang spectral]
\[
\text{rang\_cible}(k,n) = \text{rang\_base}(k) + n - 10,
\qquad
P(k,n) = \text{premier}(\text{rang\_cible}),
\]
\[
\text{Digamma}(k,n) = \SB(k,n) - P(k,n) \times k^{6}.
\]
\end{boiteformule}

\subsection{Vérification $k = 3$ à $9$ (tous corrects, 100\,\%)}

\begin{center}
\small
\begin{tabular}{@{}c r c r@{}}
\toprule
\rowcolor{spectralgris}
$k$ & $P(n=10)$ & Branche HOL & Calculé \checkmark \\
\midrule
3 & 227 & pos8 $-$ & 227 \\
4 & 947 & pos8 $+$ & 947 \\
5 & 2\,999 & pos7 $+$ & 2\,999 \\
6 & 7\,529 & pos8 $+$ & 7\,529 \\
7 & 16\,519 & pos8 $-$ & 16\,519 \\
8 & 32\,327 & pos8 $-$ & 32\,327 \\
9 & 58\,337 & pos7 $-$ & 58\,337 \\
\bottomrule
\end{tabular}
\end{center}

\subsection{Tableau complet $P(1/k)$, $n = 10$ --- $k = 10$ à $k = 111$}

Légende : \textbf{P unique} (un seul candidat standard premier) $\mid$
\textbf{C\_MULTIPLE} (plusieurs candidats premiers ; aucune sélection
automatique) $\mid$ \textbf{AUCUN P} (aucun candidat standard premier ; une
règle spéciale requiert une validation distincte). Lorsqu'une valeur et une
branche sont indiquées pour un cas multiple, elles désignent le choix de
référence documenté, pas l'unique candidat.

{\small
\begin{longtable}{@{}c r l c@{}}
\toprule
\rowcolor{spectralgris}
$k$ & $P(1/k,\ n=10)$ & État & Branche \\
\midrule
\endfirsthead
\toprule
\rowcolor{spectralgris}
$k$ & $P(1/k,\ n=10)$ & État & Branche \\
\midrule
\endhead
10 & 99\,109 & C\_MULTIPLE & pos8 $-$ \\
11 & --- & RÈGLE SPÉCIALE (PDF : 1\,611\,851) & --- \\
12 & 247\,259 & & pos8 $-$ \\
13 & 368\,939 & & pos8 $+$ \\
14 & --- & AUCUN P & --- \\
15 & 755\,789 & & pos8 $+$ \\
16 & 1\,044\,751 & C\_MULTIPLE & pos8 $-$ \\
17 & 1\,414\,943 & & pos7 $+$ \\
18 & 1\,883\,429 & & pos8 $+$ \\
19 & 2\,469\,277 & & pos7 $-$ \\
20 & 3\,192\,419 & C\_MULTIPLE & pos8 $-$ \\
21 & 4\,074\,419 & C\_MULTIPLE & pos8 $+$ \\
22 & --- & AUCUN P & --- \\
23 & 6\,424\,727 & & pos8 $-$ \\
24 & 7\,949\,399 & C\_MULTIPLE & pos8 $-$ \\
25 & 9\,750\,649 & C\_MULTIPLE & pos8 $-$ \\
26 & 11\,863\,799 & & pos7 $+$ \\
27 & --- & RÈGLE SPÉCIALE HOL ($P = 14\,330\,707$) & --- \\
28 & --- & AUCUN P & --- \\
29 & --- & AUCUN P & --- \\
30 & 24\,273\,929 & & pos8 $-$ \\
31 & --- & AUCUN P & --- \\
32 & 33\,522\,719 & & pos8 $-$ \\
33 & --- & AUCUN P & --- \\
34 & --- & AUCUN P & --- \\
35 & --- & AUCUN P & --- \\
36 & 60\,420\,851 & C\_MULTIPLE & pos8 $-$ \\
37 & 69\,293\,303 & & pos7 $+$ \\
38 & --- & AUCUN P & --- \\
39 & --- & AUCUN P & --- \\
40 & 102\,337\,639 & C\_MULTIPLE & pos8 $-$ \\
41 & --- & AUCUN P & --- \\
42 & 130\,617\,143 & & pos7 $+$ \\
43 & 146\,929\,021 & & pos7 $-$ \\
44 & --- & AUCUN P & --- \\
45 & --- & AUCUN P & --- \\
46 & 205\,867\,801 & C\_MULTIPLE & pos8 $-$ \\
47 & --- & AUCUN P & --- \\
48 & --- & AUCUN P & --- \\
49 & 282\,357\,599 & C\_MULTIPLE & pos7 $+$ \\
50 & --- & AUCUN P & --- \\
51 & --- & AUCUN P & --- \\
52 & 380\,066\,179 & & pos8 $-$ \\
53 & --- & AUCUN P & --- \\
54 & 459\,010\,529 & C\_MULTIPLE & pos8 $-$ \\
55 & --- & AUCUN P & --- \\
56 & --- & AUCUN P & --- \\
57 & 601\,506\,977 & & pos7 $-$ \\
58 & 656\,165\,077 & & pos8 $-$ \\
59 & --- & AUCUN P & --- \\
60 & --- & AUCUN P & --- \\
61 & --- & AUCUN P & --- \\
62 & 915\,894\,503 & & pos7 $+$ \\
63 & 992\,182\,589 & C\_MULTIPLE & pos8 $+$ \\
64 & 1\,073\,483\,839 & C\_MULTIPLE (3 premiers !) & pos8 $-$ \\
65 & 1\,160\,020\,289 & C\_MULTIPLE & pos8 $-$ \\
66 & 1\,252\,049\,501 & C\_MULTIPLE & pos8 $-$ \\
67 & --- & AUCUN P & --- \\
68 & --- & AUCUN P & --- \\
69 & 1\,563\,702\,839 & & pos7 $+$ \\
70 & 1\,680\,356\,999 & & pos7 $+$ \\
71 & 1\,803\,876\,551 & & pos8 $-$ \\
72 & --- & AUCUN P & --- \\
73 & --- & AUCUN P & --- \\
74 & --- & AUCUN P & --- \\
75 & 2\,372\,624\,999 & & pos7 $+$ \\
76 & --- & AUCUN P & --- \\
77 & 2\,706\,333\,629 & & pos8 $-$ \\
78 & 2\,886\,705\,977 & & pos8 $-$ \\
79 & --- & AUCUN P & --- \\
80 & 3\,276\,294\,479 & & pos8 $-$ \\
81 & --- & AUCUN P & --- \\
82 & --- & AUCUN P & --- \\
83 & --- & AUCUN P & --- \\
84 & --- & AUCUN P & --- \\
85 & 4\,436\,438\,999 & & pos7 $+$ \\
86 & --- & AUCUN P & --- \\
87 & 4\,983\,550\,877 & & pos7 $-$ \\
88 & 5\,276\,645\,527 & C\_MULTIPLE & pos8 $-$ \\
89 & --- & AUCUN P & --- \\
90 & --- & AUCUN P & --- \\
91 & --- & AUCUN P & --- \\
92 & --- & AUCUN P & --- \\
93 & --- & AUCUN P & --- \\
94 & --- & AUCUN P & --- \\
95 & --- & AUCUN P & --- \\
96 & 8\,152\,842\,431 & & pos7 $-$ \\
97 & 8\,586\,418\,271 & C\_MULTIPLE & pos8 $+$ \\
98 & --- & AUCUN P & --- \\
99 & --- & AUCUN P & --- \\
100 & --- & AUCUN P & --- \\
101 & --- & AUCUN P & --- \\
102 & 11\,039\,747\,027 & & pos7 $-$ \\
103 & 11\,591\,637\,509 & & pos8 $+$ \\
104 & --- & AUCUN P & --- \\
105 & 12\,761\,657\,999 & & pos7 $+$ \\
106 & 13\,381\,076\,101 & & pos8 $-$ \\
107 & 14\,024\,303\,819 & & pos8 $-$ \\
108 & 14\,692\,021\,271 & & pos7 $-$ \\
109 & 15\,384\,932\,747 & & pos8 $+$ \\
110 & --- & AUCUN P & --- \\
111 & 16\,849\,226\,351 & C\_MULTIPLE & pos8 $-$ \\
\bottomrule
\end{longtable}}

\subsection{Synthèse statistique ($k = 10$ à $111$)}

\begin{center}
\small
\begin{tabular}{@{}l c c@{}}
\toprule
\rowcolor{spectralgris}
État & Nombre de $k$ & \% \\
\midrule
P certifié unique & 38 & 37\,\% \\
C\_MULTIPLE ($\ge 2$ premiers) & 20 & 20\,\% \\
AUCUN P (règle standard) & 44 & 43\,\% \\
\bottomrule
\end{tabular}
\end{center}

\subsection{Notes importantes pour la programmation de l'agent}

\begin{boitealerte}[Cas AUCUN P (44 valeurs)]
Les 4 candidats standards sont tous composites. Deux stratégies possibles selon
votre HOL :
\begin{itemize}[nosep]
  \item Règle spéciale cataloguée (comme $k=27$ : Digamma $= \SA - (2k^{8}-k^{6})$, $P = 14\,330\,707$) ;
  \item Mode \textsc{création d'ancrage} : l'agent retourne C\_NON\_CERTIFIÉ et attend validation HOL externe.
\end{itemize}
\end{boitealerte}

\begin{boitenote}[Cas C\_MULTIPLE (20 valeurs)]
Selon le contrat HOL, retourner toutes les branches premières et ne faire
aucune sélection silencieuse. Une branche de référence portée par le catalogue
doit être identifiée comme telle ; elle ne rend pas les autres branches
inexistantes ni le cas unique.
\end{boitenote}

\begin{boitenote}[Règles spéciales]
$k=11$ ne donne aucun candidat premier aux quatre essais standards. La valeur
$P=1\,611\,851$ (rang 121\,982), associée à la règle A7$-$ du moteur géométrique,
est une règle spéciale distincte et doit être validée séparément. De même,
$k=27$ utilise une règle spéciale HOL documentée ; aucune de ces règles ne doit
être extrapolée à d'autres rapports.
\end{boitenote}

\begin{boitealerte}[Arithmétique exacte requise]
Pour $k \ge 50$, les valeurs de $\SA$ et $\SB$ dépassent $10^{15}$ : l'agent
doit utiliser l'arithmétique entière Python (pas float) pour éviter les erreurs
d'arrondi Excel.
\end{boitealerte}

\subsection{Exécutions multi-rapports archivées ($n=10$)}

Deux rapports d'exécution du dépôt couvrent les plages communes suivantes :
\url{liste_k2_k1001_n10.txt} traite $k=2..1001$ (1\,000 rapports) et
annonce 423 succès ; \url{liste_k1002_k10001_n10.txt} traite
$k=1002..10001$ (9\,000 rapports) et annonce 2\,604 succès. Ces compteurs
reproduisent le statut du dispatcher v7.9 au moment de la génération.

\begin{boitealerte}[Portée des rapports archivés]
Ces fichiers sont des journaux de calcul, pas des certificats indépendants.
Le dispatcher historique retenait le premier candidat premier rencontré ; il
ne consignait pas systématiquement toutes les branches premières. Les succès
ne démontrent donc ni l'unicité de l'ancre, ni la primalité démontrée par HOL,
ni une propriété sur les zéros de la fonction zêta. Le dispatcher corrigé
conserve les branches candidates, signale les ambiguïtés hors catalogue et
réserve les sélections documentées à une décision explicite.
\end{boitealerte}

% ======================================================================
\section{Exemple numérique complet --- $n = 15$}
% ======================================================================

\subsection{Rapport $1/13$ ($k = 13$, $n = 15$)}

\subsubsection{Niveau 1 --- Constantes Savard universelles}

\begin{center}\small
\begin{tabular}{@{}l l l@{}}
\toprule
\rowcolor{spectralgris}
Constante & Formule & Valeur exacte \\
\midrule
$\alpha_A(13)$ & $2(13^{4}-13^{2}+1)/(12\cdot 13^{3})$ & $28\,393/13\,182$ \\
$\alpha_B(13)$ & $13\cdot\alpha_A(13)$ & $28\,393/1014$ \\
$\delta_A(13)$ & $13/12$ & $13/12$ \\
$\delta_B(13)$ & $(13^{7}-13^{6}+13)/12$ & $57\,921\,721/12$ \\
\bottomrule
\end{tabular}
\end{center}

\subsubsection{Niveau 2 --- Formules fermées}
\begin{align*}
\SA(13,15) &= \frac{(k^{4}-k^{2}+1)k^{12}}{k-1}-\frac{k}{k-1},\\
&= \frac{28\,393\times23\,298\,085\,122\,481}{12}-\frac{13}{12},\\
&=55\,125\,210\,906\,883\,585,\\
\SB(13,15) &= \frac{(k^{4}-k^{2}+1)k^{13}}{k-1}
-\frac{k^{7}-k^{6}+k}{k-1},\\
&=716\,627\,741\,784\,659\,809.
\end{align*}

Invariant spectral (Théorème XIV.3) :
\begin{align*}
\Delta\SA &= \SA(n{=}15)-\SA(n{=}10),\\
&=55\,125\,210\,906\,883\,585-148\,468\,220\,264,\\
&=55\,125\,062\,438\,663\,321,\\
\Delta\SB &= \SB(n{=}15)-\SB(n{=}10),\\
&=716\,627\,741\,784\,659\,809-1\,930\,082\,036\,636,\\
&=716\,625\,811\,702\,623\,173,\\
\Delta\SA / \Delta\SB &= 1/13 \quad \checkmark
\end{align*}

\subsubsection{Niveau 3 --- Construction terme à terme ($n = 15$)}

\begin{center}\small
\begin{tabular}{@{}c r r l l@{}}
\toprule
\rowcolor{spectralgris}
$i$ & $A(13,15,i)$ & $B(13,15,i)$ & Règle A & Règle B \\
\midrule
1 & 13 & 13 & $k^{1}$ & $k^{1}$ \\
2 & 169 & 169 & $k^{2}$ & $k^{2}$ \\
3 & 2\,197 & 2\,197 & $k^{3}$ & $k^{3}$ \\
4 & 28\,561 & 28\,561 & $k^{4}$ & $k^{4}$ \\
5 & 371\,293 & 371\,293 & $k^{5}$ & $k^{5}$ \\
6 & 4\,826\,809 & 62\,748\,517 & $k^{6}$ & $k^{7}$ $\leftarrow$ Saut Zêta \\
7 & 62\,748\,517 & 815\,730\,721 & $k^{7}$ & $k^{8}$ \\
8 & 815\,730\,721 & 10\,604\,499\,373 & $k^{8}$ & $k^{9}$ \\
9 & 10\,604\,499\,373 & 137\,858\,491\,849 & $k^{9}$ & $k^{10}$ \\
10 & 137\,858\,491\,849 & 1\,792\,160\,394\,037 & $k^{10}$ & $k^{11}$ \\
11 & 1\,792\,160\,394\,037 & 23\,298\,085\,122\,481 & $k^{11}$ & $k^{12}$ \\
12 & 23\,298\,085\,122\,481 & 302\,875\,106\,592\,253 & $k^{12}$ & $k^{13}$ \\
13 & 302\,875\,106\,592\,253 & 3\,937\,376\,385\,699\,289 & $k^{13}$ & $k^{14}$ \\
14 & 3\,914\,078\,300\,576\,808 & 50\,883\,017\,907\,498\,504 & $k^{14}-k^{12}$ & $k^{15}-k^{13}$ \\
15 & 50\,883\,017\,907\,498\,504 & 661\,479\,232\,797\,480\,552 & $k^{15}-k^{13}$ & $k^{16}-k^{14}$ \\
\midrule
$\Sigma$ & \multicolumn{2}{l}{$\Sigma A = 55\,125\,210\,906\,883\,585 = \SA(13,15)$} & \multicolumn{2}{l}{$\Sigma B = 716\,627\,741\,784\,659\,809 = \SB(13,15)$} \\
\bottomrule
\end{tabular}
\end{center}

\subsubsection{Niveau 4 --- 4 essais Digamma ($k^{6} = 4\,826\,809$)}

Les 4 essais sont calibrés pour $n=10$. Pour $n=15$, ils servent de
vérification structurelle --- tous composites est le comportement attendu.

\begin{center}\small
\begin{tabular}{@{}c r r l@{}}
\toprule
\rowcolor{spectralgris}
Branche & Digamma & $P$ candidat & État \\
\midrule
pos13$-$ & 54\,822\,335\,800\,291\,332 & 137\,110\,336\,453 & composé \\
pos13$+$ & 55\,428\,086\,013\,475\,838 & 136\,984\,839\,419 & composé \\
pos12$+$ & 55\,148\,508\,992\,006\,066 & 137\,042\,761\,127 & composé \\
pos12$-$ & 55\,101\,912\,821\,761\,104 & 137\,052\,414\,745 & composé \\
\bottomrule
\end{tabular}
\end{center}

\subsubsection{Niveau 5 --- Reconstruction rang cible (méthode officielle pour $n \neq 10$)}
\begin{gather*}
P(13,n{=}10)\ \text{documenté}=368\,939,\\
\text{rang\_base}(13)=\pi(368\,939)=31\,452,\\
\Delta n=15-10=5,\\
\text{rang\_cible}(13,n{=}15)=31\,452+5=31\,457,\\
P(13,n{=}15)=\text{premier}(31\,457)=369\,023
\quad\textbf{PREMIER certifié}.
\end{gather*}
\begin{multline*}
\text{Digamma}(13,15)=\SB-P\times k^{6}\\
=716\,627\,741\,784\,659\,809\\
{}-369\,023\times4\,826\,809\\
=716\,625\,960\,581\,122\,202.
\end{multline*}

\subsection{Rapport $1/42$ ($k = 42$, $n = 15$)}

\subsubsection{Niveau 1 --- Constantes Savard universelles}

\begin{center}\small
\begin{tabular}{@{}l l@{}}
\toprule
\rowcolor{spectralgris}
Constante & Valeur exacte \\
\midrule
$\alpha_A(42)$ & $3\,109\,933\,/\,1\,518\,804$ \\
$\alpha_B(42)$ & $3\,109\,933\,/\,36\,162$ \\
$\delta_A(42)$ & $42/41$ \\
$\delta_B(42)$ & $225\,050\,301\,546\,/\,41$ \\
\bottomrule
\end{tabular}
\end{center}

\subsubsection{Niveau 2 --- Formules fermées}
\begin{align*}
\SA(42,15)&=2\,285\,381\,254\,365\,702\,554\,104\,806,\\
\SB(42,15)&=95\,986\,012\,683\,359\,501\,783\,370\,150.
\end{align*}

\subsubsection{Niveau 3 --- Construction terme à terme ($n = 15$)}

\begin{center}\small
\begin{tabular}{@{}c r r@{}}
\toprule
\rowcolor{spectralgris}
$i$ & $A(42,15,i)$ & $B(42,15,i)$ \\
\midrule
1 & 42 & 42 \\
2 & 1\,764 & 1\,764 \\
3 & 74\,088 & 74\,088 \\
4 & 3\,111\,696 & 3\,111\,696 \\
5 & 130\,691\,232 & 130\,691\,232 \\
6 & 5\,489\,031\,744 & 230\,539\,333\,248 \quad $\leftarrow$ Saut Zêta ($k^{7}$) \\
7 & 230\,539\,333\,248 & 9\,682\,651\,996\,416 \\
8 & 9\,682\,651\,996\,416 & 406\,671\,383\,849\,472 \\
9 & 406\,671\,383\,849\,472 & 17\,080\,198\,121\,677\,824 \\
10 & 17\,080\,198\,121\,677\,824 & 717\,368\,321\,110\,468\,608 \\
11 & 717\,368\,321\,110\,468\,608 & 30\,129\,469\,486\,639\,681\,536 \\
12 & 30\,129\,469\,486\,639\,681\,536 & 1\,265\,437\,718\,438\,866\,624\,512 \\
13 & 1\,265\,437\,718\,438\,866\,624\,512 & 53\,148\,384\,174\,432\,398\,229\,504 \\
14 & 53\,118\,254\,704\,945\,758\,547\,968 & 2\,230\,966\,697\,607\,721\,859\,014\,656 \\
15 & 2\,230\,966\,697\,607\,721\,859\,014\,656 & 93\,700\,601\,299\,524\,318\,078\,615\,552 \\
\midrule
$\Sigma$ & \multicolumn{2}{l}{$\Sigma A = 2\,285\,381\,254\,365\,702\,554\,104\,806$ \quad $\Sigma B = 95\,986\,012\,683\,359\,501\,783\,370\,150$} \\
\bottomrule
\end{tabular}
\end{center}

\subsubsection{Niveau 4 --- 4 essais Digamma ($k^{6} = 5\,489\,031\,744$)}

\begin{center}\small
\begin{tabular}{@{}c r r l@{}}
\toprule
\rowcolor{spectralgris}
Branche & Digamma & $P$ candidat & État \\
\midrule
pos13$-$ & 2\,284\,115\,816\,647\,263\,687\,480\,294 & 17\,070\,751\,498\,046\,399 & composé \\
pos13$+$ & 2\,286\,646\,692\,084\,141\,420\,729\,318 & 17\,070\,290\,419\,379\,903 & composé \\
pos12$+$ & 2\,285\,411\,383\,835\,189\,193\,786\,342 & 17\,070\,515\,469\,681\,407 & composé \\
pos12$-$ & 2\,285\,351\,124\,896\,215\,914\,423\,270 & 17\,070\,526\,447\,744\,895 & composé \\
\bottomrule
\end{tabular}
\end{center}

\subsubsection{Niveau 5 --- Reconstruction rang\_cible}
\begin{gather*}
P(42,n{=}10)\ \text{documenté}=130\,617\,143,\\
\text{rang\_base}(42)=\pi(130\,617\,143)=7\,411\,117,\\
\Delta n=15-10=5,\\
\text{rang\_cible}(42,n{=}15)=7\,411\,117+5=7\,411\,122,\\
P(42,n{=}15)=\text{premier}(7\,411\,122)=130\,617\,251
\quad\textbf{PREMIER certifié}.
\end{gather*}
\begin{multline*}
\text{Digamma}(42,15)=\SB-P\times k^{6}\\
=95\,986\,012\,683\,359\,501\,783\,370\,150\\
{}-130\,617\,251\times5\,489\,031\,744\\
=95\,986\,011\,966\,397\,264\,730\,354\,406.
\end{multline*}

\subsection{Rapport $1/54$ ($k = 54$, $n = 15$)}

\subsubsection{Niveau 1 --- Constantes Savard universelles}

\begin{center}\small
\begin{tabular}{@{}l l@{}}
\toprule
\rowcolor{spectralgris}
Constante & Valeur exacte \\
\midrule
$\alpha_A(54)$ & $8\,500\,141\,/\,4\,172\,796$ \\
$\alpha_B(54)$ & $8\,500\,141\,/\,77\,274$ \\
$\delta_A(54)$ & $54/53$ \\
$\delta_B(54)$ & $1\,314\,130\,298\,742\,/\,53$ \\
\bottomrule
\end{tabular}
\end{center}

\subsubsection{Niveau 2 --- Formules fermées}
\begin{align*}
\SA(54,15)&=98\,599\,651\,085\,954\,948\,762\,648\,034,\\
\SB(54,15)&=5\,324\,381\,158\,641\,567\,208\,388\,082\,594,\\
\Delta\SA/\Delta\SB&=1/54\quad\checkmark
\end{align*}

\subsubsection{Niveau 3 --- Construction terme à terme ($n = 15$)}

\begin{center}\small
\begin{tabular}{@{}c r r@{}}
\toprule
\rowcolor{spectralgris}
$i$ & $A(54,15,i)$ & $B(54,15,i)$ \\
\midrule
1--5 & \multicolumn{2}{l}{54 / 2\,916 / 157\,464 / 8\,503\,056 / 459\,165\,024 \quad (idem A et B)} \\
6 & 24\,794\,911\,296 & 1\,338\,925\,209\,984 \quad $\leftarrow$ Saut Zêta ($54^{7}$) \\
7 & 1\,338\,925\,209\,984 & 72\,301\,961\,339\,136 \\
8 & 72\,301\,961\,339\,136 & 3\,904\,305\,912\,313\,344 \\
9 & 3\,904\,305\,912\,313\,344 & 210\,832\,519\,264\,920\,576 \\
10 & 210\,832\,519\,264\,920\,576 & 11\,384\,956\,040\,305\,711\,104 \\
11 & 11\,384\,956\,040\,305\,711\,104 & 614\,787\,626\,176\,508\,399\,616 \\
12 & 614\,787\,626\,176\,508\,399\,616 & 33\,198\,531\,813\,531\,453\,579\,264 \\
13 & 33\,198\,531\,813\,531\,453\,579\,264 & 1\,792\,720\,717\,930\,698\,493\,280\,256 \\
14 & 1\,792\,105\,930\,304\,521\,984\,880\,640 & 96\,773\,720\,236\,444\,187\,183\,554\,560 \\
15 & 96\,773\,720\,236\,444\,187\,183\,554\,560 & 5\,225\,780\,892\,767\,986\,107\,911\,946\,240 \\
\midrule
$\Sigma$ & \multicolumn{2}{l}{$\Sigma A = 98\,599\,651\,085\,954\,948\,762\,648\,034$ \quad $\Sigma B = 5\,324\,381\,158\,641\,567\,208\,388\,082\,594$} \\
\bottomrule
\end{tabular}
\end{center}

\subsubsection{Niveau 4 --- 4 essais Digamma ($k^{6} = 24\,794\,911\,296$)}

\begin{center}\small
\begin{tabular}{@{}c >{\raggedleft\arraybackslash}p{5.6cm} >{\raggedleft\arraybackslash}p{4.2cm} l@{}}
\toprule
\rowcolor{spectralgris}
Branche & Digamma & $P$ candidat & État \\
\midrule
pos13$-$ & 98\,566\,452\,554\,141\,417\,309\,068\,770 & 210\,761\,581\,023\,702\,719 & composé \\
pos13$+$ & 98\,632\,849\,617\,768\,480\,216\,227\,298 & 210\,758\,903\,173\,282\,751 & composé \\
pos12$+$ & 98\,600\,265\,873\,581\,125\,271\,047\,650 & 210\,760\,217\,303\,581\,439 & composé \\
pos12$-$ & 98\,599\,036\,298\,328\,772\,254\,248\,418 & 210\,760\,266\,893\,404\,031 & composé \\
\bottomrule
\end{tabular}
\end{center}

\subsubsection{Niveau 5 --- Reconstruction rang\_cible}
\begin{gather*}
P(54,n{=}10)\ \text{documenté}=459\,010\,529,\\
\text{rang\_base}(54)=\pi(459\,010\,529)=24\,304\,918,\\
\Delta n=15-10=5,\\
\text{rang\_cible}(54,n{=}15)=24\,304\,918+5=24\,304\,923,\\
P(54,n{=}15)=\text{premier}(24\,304\,923)=459\,010\,597
\quad\textbf{PREMIER certifié}.
\end{gather*}
\begin{multline*}
\text{Digamma}(54,15)=\SB-P\times k^{6}\\
=5\,324\,381\,158\,641\,567\,208\,388\,082\,594\\
{}-459\,010\,597\times24\,794\,911\,296\\
=5\,324\,381\,147\,260\,440\,171\,849\,078\,882.
\end{multline*}
\begin{boitealerte}[Statut de l'ancre $1/54$]
Le tableau 3.7 classe $k=54$ comme \textbf{C\_MULTIPLE}. Les essais exacts
produisent au moins les deux branches premières $A(8)-=459\,010\,529$ et
$A(8)+=459\,004\,697$. Les calculs de rang ci-dessus suivent la branche
$A(8)-$ explicitement documentée par la source ; le résultat n'est pas une
ancre unique et doit rester conditionnel à cette sélection de référence.
\end{boitealerte}

\subsection{Tableau récapitulatif final --- $P(1/k,\ n = 15)$}

\begin{center}\small
\resizebox{\textwidth}{!}{%
\begin{tabular}{@{}c r r r r r r l@{}}
\toprule
\rowcolor{spectralgris}
$k$ & \multicolumn{1}{c}{$\SA(k,15)$} & \multicolumn{1}{c}{$\SB(k,15)$} & rang base & rang cible & $P(k,n{=}15)$ & $P(k,n{=}10)$ & État \\
\midrule
13 & 55\,125\,210\,906\,883\,585 & 716\,627\,741\,784\,659\,809 & 31\,452 & 31\,457 & 369\,023 & 368\,939 & CERTIFIÉ \\
42 & $2.29\times 10^{24}$ & $9.60\times 10^{25}$ & 7\,411\,117 & 7\,411\,122 & 130\,617\,251 & 130\,617\,143 & CERTIFIÉ \\
54 & $9.86\times 10^{25}$ & $5.32\times 10^{27}$ & 24\,304\,918 & 24\,304\,923 & 459\,010\,597 & 459\,010\,529 & C\_MULTIPLE ; A(8)- documentée \\
\bottomrule
\end{tabular}
}
\end{center}

\subsection{Notes critiques pour l'agent}

\subsubsection{1. Saut Zêta --- asymétrie A / B confirmée}
La correction importante validée ici : $A(k,n,i) = k^{i}$ pour tout
$1 \le i \le n-2$ (pas de Saut Zêta dans la suite A). Le Saut Zêta $k^{7}$ à la
position $i=6$ n'existe que dans la suite B. Les deux suites sont identiques
pour $i \le 5$, puis divergent à partir de $i = 6$.

\subsubsection{2. Niveau 4 hors-calibrage pour $n \neq 10$}
Les 4 essais Digamma (pos$(n-2)\pm$, pos$(n-3)\pm$) produisent tous des
candidats composites pour $n=15$ --- ce comportement est attendu et correct.
Ces essais sont calibrés structurellement sur $n=10$ ; le Niveau 5 est la voie
officielle pour tout $n \neq 10$.

\subsubsection{3. Niveau 5 --- formule rang\_cible}
\begin{boiteformule}[Pseudo-code officiel]
\texttt{rang\_base(k) = primepi(P(k, n=10))} \quad \# rang du P certifié à n=10\\
\texttt{rang\_cible(k, n) = rang\_base(k) + (n - 10)}\\
\texttt{P(k, n) = prime(rang\_cible(k, n))} \quad \# arithmétique exacte sympy
\end{boiteformule}
L'écart en rang $+5$ (pour $\Delta n = 5$) est cohérent pour les trois $k$,
et illustre la règle de rang pour ces trois exemples. Trois cas numériques ne
suffisent pas à établir l'universalité de cette règle pour tout $k$.

\subsubsection{4. Invariant spectral (Théorème XIV.3) --- vérifié pour les 3 rapports}
\[
\Delta\SA(k, n{=}15, n{=}10)\,/\,\Delta\SB(k, n{=}15, n{=}10) = 1/k
\qquad \forall\, k \in \{13, 42, 54\} \quad \checkmark
\]

\subsubsection{5. Arithmétique --- seuils de dépassement float}

\begin{center}\small
\begin{tabular}{@{}c c c l@{}}
\toprule
\rowcolor{spectralgris}
$k$ & $\SA(k,15)$ & Ordre de grandeur & Float 64 safe ? \\
\midrule
13 & 55\,125\,210\,906\,883\,585 & $\sim 5.5\times 10^{16}$ & limite \\
42 & $2.29\times 10^{24}$ & $\sim 10^{24}$ & erreur garantie \\
54 & $9.86\times 10^{25}$ & $\sim 10^{25}$ & erreur garantie \\
\bottomrule
\end{tabular}
\end{center}

\begin{boitealerte}
Arithmétique entière Python obligatoire dès $k \ge 42$ pour $n = 15$.
Excel / float64 introduiraient des erreurs d'arrondi fatales.
\end{boitealerte}


% ======================================================================
\section{Schéma neuronal et arborescence algorithmique du système convolutif}
% ======================================================================

Cette section est construite à partir des tables convolutives du classeur Excel
\url{systeme_convolutif_spectral_general.xlsx} (feuilles
\emph{Convolution}, \emph{Systeme General}, \emph{Parametres},
\emph{Moteur Géo Convolutif}, \emph{Tests Digamma}, \emph{Suites Géo A-B}).

\subsection{Schéma neuronal de la table convolutive}

Chaque terme de la table de convolution formelle vaut $c_1 \cdot k^{e_1} +
c_2 \cdot k^{e_2}$ : cette représentation encode les additions et soustractions
sans figer $k$. Le schéma ci-dessous en trace le réseau de propagation, de
l'entrée $(k, n)$ jusqu'au premier certifié.

\begin{figure}[h!]
\centering
\begin{tikzpicture}[
  font=\footnotesize,
  entree/.style={rectangle,rounded corners,draw=spectralbleu,fill=spectralbleu!12,thick,minimum width=2.6cm,minimum height=0.8cm,align=center},
  neurone/.style={circle,draw=spectralazur,fill=spectralazur!12,thick,minimum size=0.62cm,inner sep=1pt},
  neuroneB/.style={circle,draw=spectralviolet,fill=spectralviolet!12,thick,minimum size=0.62cm,inner sep=1pt},
  somme/.style={rectangle,rounded corners,draw=spectralor,fill=spectralor!18,thick,minimum width=1.7cm,minimum height=0.8cm,align=center},
  digamma/.style={diamond,draw=spectralviolet,fill=spectralviolet!10,thick,aspect=2.2,align=center,inner sep=0.5pt},
  sortie/.style={rectangle,rounded corners,draw=spectralbleu,fill=spectralbleu,text=white,thick,minimum width=2.9cm,minimum height=0.9cm,align=center},
  flux/.style={-{Stealth[length=2.2mm]},thick,spectralbleu!70},
  fluxor/.style={-{Stealth[length=2.2mm]},thick,spectralor!90!black}
]
  % --- Entrées
  \node[entree] (k) at (0,0) {Rapport $1/k$\\[-1pt]$k \ge 2$ entier};
  \node[entree] (n) at (3.4,0) {$n$ = nombre de\\[-1pt]termes ($n \ge 8$)};

  % --- Couche termes (neurones)
  \foreach \i/\lab in {1/$k^{1}$,2/$k^{2}$,3/$k^{3}$,5/$k^{i}$,7/$k^{n-2}$}{
    \node[neurone] (a\i) at (0.4*\i+0.2,-2.1) {\lab};}
  \node at (3.4,-2.1) {\footnotesize $\cdots$};
  \node[neurone,draw=spectralor,fill=spectralor!25] (as1) at (3.6,-2.1) {};
  \node[neurone,draw=spectralor,fill=spectralor!25] (as2) at (4.4,-2.1) {};
  \node[below=1pt of as1,xshift=-4pt] {\scriptsize Savard $k^{n-1}{-}k^{n-3}$};
  \node[below=1pt of as2,xshift=4pt] {\scriptsize Savard $k^{n}{-}k^{n-2}$};

  \foreach \i/\lab in {1/$k^{1}$,2/$k^{2}$,3/$k^{5}$,5/$k^{i+1}$,7/$k^{n-2}$}{
    \node[neuroneB] (b\i) at (0.4*\i+0.2,-3.4) {\lab};}
  \node at (3.4,-3.4) {\footnotesize $\cdots$};
  \node[neuroneB,draw=spectralor,fill=spectralor!25] (bs1) at (3.6,-3.4) {};
  \node[neuroneB,draw=spectralor,fill=spectralor!25] (bs2) at (4.4,-3.4) {};
  \node[above=1pt of b3,yshift=2pt] {};
  \node[below=1pt of bs1,xshift=-4pt] {\scriptsize décalé $k^{n}{-}k^{n-2}$};
  \node[below=1pt of bs2,xshift=4pt] {\scriptsize décalé $k^{n+1}{-}k^{n-1}$};
  \node[neuroneB,draw=spectralrouge,fill=spectralrouge!15] (zeta) at (2.6,-3.4) {$k^{7}$};
  \node[above=1pt of zeta] {\scriptsize\color{spectralrouge} Saut Zêta ($i=6$)};

  % --- Sommes
  \node[somme] (sa) at (2.1,-5.0) {$\SA(k,n)$\\[-2pt]\scriptsize $\frac{\alpha_A}{2}k^{n}-\delta_A$};
  \node[somme] (sb) at (2.1,-6.3) {$\SB(k,n)$\\[-2pt]\scriptsize $\frac{\alpha_B}{2}k^{n}-\delta_B$};

  % --- Digamma : 4 branches
  \node[digamma] (d1) at (6.9,-4.4) {\scriptsize $\SA - k^{8}$};
  \node[digamma] (d2) at (6.9,-5.3) {\scriptsize $\SA + k^{8}$};
  \node[digamma] (d3) at (6.9,-6.2) {\scriptsize $\SA + k^{7}$};
  \node[digamma] (d4) at (6.9,-7.1) {\scriptsize $\SA - k^{7}$};

  % --- Division Zêta + primalité
  \node[somme,minimum width=2.5cm] (div) at (10.6,-5.3) {$P = \dfrac{\SB - \text{Digamma}}{k^{6}}$\\[-1pt]\scriptsize test : entier ? premier ?};
  \node[sortie] (p) at (10.6,-7.4) {$P$ certifié\\\scriptsize ancrage $n=10$};
  \node[sortie,fill=spectralviolet] (rang) at (10.6,-8.9) {Niveau 5\\\scriptsize rang cible $=$ rang base $+\,n-10$};

  % --- Flux
  \draw[flux] (k.south) -- (1.1,-1.2) -| (a1.north);
  \draw[flux] (n.south) -- (3.4,-1.2) -| (as2.north);
  \foreach \i in {1,2,3,5,7}{ \draw[fluxor] (a\i.south) -- (a\i.south |- sa.north); }
  \draw[fluxor] (as1.south) -- (as1.south |- sa.north);
  \draw[fluxor] (as2.south) -- (as2.south |- sa.north);
  \foreach \i in {1,2,3,5,7}{ \draw[fluxor] (b\i.south) -- (b\i.south |- sb.north); }
  \draw[fluxor] (zeta.south) -- (zeta.south |- sb.north);
  \draw[fluxor] (bs1.south) -- (bs1.south |- sb.north);
  \draw[fluxor] (bs2.south) -- (bs2.south |- sb.north);
  \draw[flux] (sa.east) -- (d1.west); \draw[flux] (sa.east) -- (d2.west);
  \draw[flux] (sa.east) -- (d3.west); \draw[flux] (sa.east) -- (d4.west);
  \draw[flux] (sb.east) -| (div.south);
  \foreach \i in {1,2,3,4}{ \draw[flux] (d\i.east) -- (div.west |- d\i.east); }
  \draw[flux] (div.south) -- (p.north);
  \draw[flux] (p.south) -- (rang.north) node[midway,right]{\scriptsize $n \neq 10$};
  % Étiquettes de couches
  \node[rotate=90,anchor=south,color=spectralbleu] at (-1.15,-2.75) {\scriptsize couche termes (table Convolution)};
  \node[rotate=90,anchor=south,color=spectralor!80!black] at (-1.15,-5.65) {\scriptsize couche sommes};
  \node[rotate=90,anchor=south,color=spectralviolet] at (5.6,-5.75) {\scriptsize 4 branches Digamma};
\end{tikzpicture}
\caption{Schéma neuronal du système convolutif : propagation des termes
$c_1 k^{e_1} + c_2 k^{e_2}$ des suites A et B vers les sommes $\SA$/$\SB$, puis
vers les quatre branches Digamma et la division Zêta $k^{6}$ menant au premier
certifié (ancrage $n=10$) ou à la loi de rang spectral (Niveau 5).}
\end{figure}

\subsection{Arborescence algorithmique --- niveaux 1 et 2}

Les deux niveaux évaluent la même construction sous deux représentations : le
niveau géométrique multiplie les termes entiers par un facteur commun et
reproduit les mêmes quatre candidats. Il ne crée pas une ancre supplémentaire
lorsque le niveau entier n'en trouve aucune. Le résultat à $n=10$ doit donc
être classé comme candidat unique, choix de catalogue explicite, ambiguïté sans
décision ou absence de candidat standard.

\begin{figure}[h!]
\centering
\begin{tikzpicture}[
  scale=0.82,transform shape,
  font=\footnotesize,
  etape/.style={rectangle,rounded corners,draw=spectralbleu,fill=spectralgris,thick,minimum width=6.4cm,minimum height=0.78cm,align=center},
  decision/.style={diamond,draw=spectralor,fill=spectralor!15,thick,aspect=2.6,align=center,inner sep=0.5pt},
  fin/.style={rectangle,rounded corners,draw=spectralbleu,fill=spectralbleu,text=white,thick,minimum width=4.4cm,minimum height=0.8cm,align=center},
  geo/.style={rectangle,rounded corners,draw=spectralviolet,fill=spectralviolet!10,thick,minimum width=5.4cm,minimum height=0.78cm,align=center},
  fl/.style={-{Stealth[length=2.2mm]},thick,spectralbleu!80}
]
  \node[etape] (e1) at (0,0) {\textbf{[1]} Sommes A/B à $n=10$ --- suites construites terme à terme};
  \node[etape] (e2) at (0,-1.1) {\textbf{[2]} Quatre possibilités Digamma : $\SA \pm k^{8}$, $\SA \pm k^{7}$};
  \node[etape] (e3) at (0,-2.2) {\textbf{[3]} Sommes A/B à $n=9$ (premier précédent)};
  \node[etape] (e4) at (0,-3.3) {\textbf{[4]} Coefficients : $\bigl(\mathrm{S}(10)-\mathrm{S}(9)\bigr)/k^{8}$};
  \node[etape] (e5) at (0,-4.4) {\textbf{[5]} $(\text{Reste}+x)$ $\to$ blocs A (entière) / B (décimale)};
  \node[etape] (e6) at (0,-5.5) {\textbf{[6]} Équations généralisées $\frac{\text{Coeff}}{x}\,k^{n} - \text{Reste}$ ($n > 0$ et $n < 0$)};
  \node[etape] (e7) at (0,-6.6) {\textbf{[7]} Sommes A/B au $n$ demandé $+$ premier conséquent};

  \node[decision] (d10) at (-5.6,-1.1) {$n = 10$ ?};
  \node[decision] (dpr) at (-5.6,-3.3) {$\ge 1$ premier\\parmi 4 essais ?};
  \node[geo] (g2) at (-5.6,-5.6) {\textbf{Niveau 2 (fallback géométrique)}\\suites $\sqrt{a^{2}+b^{2}}$, facteur $g(k)=\frac{\sqrt{1+k^{2}}}{k}$\\extraction floor/ceil/round, test de primalité};
  \node[fin] (f1) at (-5.6,-7.6) {$P$ certifié --- ancrage $n=10$};

  \node[decision] (dn) at (5.9,-1.1) {$n \neq 10$ ?};
  \node[etape,minimum width=5.6cm] (r1) at (5.9,-2.9) {\textbf{Niveau 5} : rang cible $=$ rang base $+\,n-10$\\$P(k,n) = \text{premier}(\text{rang cible})$};
  \node[etape,minimum width=5.6cm] (r2) at (5.9,-4.3) {Digamma à rebours : $\SB(k,n) - P(k,n)\,k^{6}$};
  \node[fin] (f2) at (5.9,-7.6) {Premier conséquent certifié};

  \draw[fl] (e1) -- (e2); \draw[fl] (e2) -- (e3); \draw[fl] (e3) -- (e4);
  \draw[fl] (e4) -- (e5); \draw[fl] (e5) -- (e6); \draw[fl] (e6) -- (e7);
  \draw[fl] (d10) -- node[above]{\scriptsize oui} (e2);
  \draw[fl] (d10.north) |- node[above,pos=0.75]{\scriptsize non} (dn.west);
  \draw[fl] (e2.west) -- (dpr.north);
  \draw[fl] (dpr) -- node[above]{\scriptsize non} (g2.north);
  \draw[fl] (dpr.east) -- node[above]{\scriptsize oui} (e3.west);
  \draw[fl] (g2) -- (f1) node[midway,right]{\scriptsize $\ge 1$ premier};
  \draw[fl] (e7.west) -| (f1.east);
  \draw[fl] (dn) -- node[right]{\scriptsize oui} (r1);
  \draw[fl] (r1) -- (r2);
  \draw[fl] (r2) -- (f2);
  \draw[fl,dashed,spectralor] (f1.south) -- ++(0,-0.5) -| (r1.south)
    node[pos=0.25,below]{\scriptsize ancrage obligatoire};
\end{tikzpicture}
\caption{Arborescence algorithmique des étapes du système convolutif pour les
niveaux 1 (entier) et 2 (géométrique), avec la voie officielle du Niveau 5 pour
tout $n \neq 10$. Invariants : le Digamma (4 possibilités) et les équations de
reconstruction pour $\forall n$ demeurent identiques d'un niveau à l'autre.}
\end{figure}

\subsection{Équations de reconstruction des suites A et B --- les deux niveaux}

\begin{boiteformule}[Niveau 1 (entier) --- reconstruction généralisée de $n$]
\begin{align*}
\SA(k,n) &= \frac{\alpha_A(k)}{2}\,k^{n} - \delta_A(k)\\
&= \frac{k^{4}-k^{2}+1}{(k-1)\,k^{3}}\,k^{n} - \frac{k}{k-1},\\
\SB(k,n) &= \frac{\alpha_B(k)}{2}\,k^{n} - \delta_B(k)\\
&= \frac{k^{4}-k^{2}+1}{(k-1)\,k^{2}}\,k^{n}\\
&\quad-\frac{k^{7}-k^{6}+k}{k-1}.
\end{align*}
formes fermées exactes valables pour tout $k$ entier $\ge 3$ et tout $n$
entier (positif ou négatif), coïncidant avec la construction terme à terme
(Niveau 3) dès $n = 8$ ; pour $n \le 7$, la construction terme à terme
progression simple fait foi (seuil des formes fermées : $n = 8$).
\end{boiteformule}

\begin{boiteformule}[Niveau 2 (géométrique) --- modèle hypoténuse $\sqrt{a^{2}+b^{2}}$]
Chaque terme géométrique est le terme entier correspondant multiplié par le
zêta géométrique $g(k) = \dfrac{\sqrt{1+k^{2}}}{k}$ :
\begin{align*}
\SA^{\,\mathrm{g\acute eo}}(k,n) &= g(k)\cdot\SA(k,n),\\
\SB^{\,\mathrm{g\acute eo}}(k,n) &= g(k)\cdot\SB(k,n),\\
\mathrm{Digamma}_{\mathrm{g\acute eo}} &= g(k)\cdot\mathrm{Digamma},\\
\mathrm{Z\hat eta}_{\mathrm{g\acute eo}} &= g(k)\cdot k^{6}.
\end{align*}
\textbf{Identité centrale :} le facteur $g(k)$ s'annule dans
\begin{align*}
P &= \frac{\SB^{\mathrm{g\acute eo}}-\mathrm{Digamma}_{\mathrm{g\acute eo}}}
{\mathrm{Z\hat eta}_{\mathrm{g\acute eo}}}\\
&= \frac{\SB-\mathrm{Digamma}}{k^{6}}.
\end{align*}
\emph{Les deux niveaux ont les mêmes utilités et les mêmes ancrages}
(8/8 conformes, Section XIV.5).
\end{boiteformule}

\begin{boitenote}[Exemple numérique croisé --- rapport $1/3$]
Au niveau 1 :
\[
\SA(3,10)=\frac{73}{9}\cdot6\cdot3^{10}-\frac{3}{2}=79\,824,
\qquad
\text{Digamma}=\SA-3^{8}=73\,263.
\]
La reconstruction donne
\[
P=\frac{238\,746-73\,263}{729}=227
\quad(49\textsuperscript{ème} premier).
\]
Au niveau 2, les sommes géométriques sont multipliées par
$g(3)=\sqrt{10}/3\approx1.05409255338946$ ; le facteur s'élimine et redonne
$P=227$.
\end{boitenote}

% ======================================================================
\section{Fiche technique --- Pipeline cognitif de l'agent Mm.\ Gabriel multiloop}
% ======================================================================

La présente fiche technique décrit le fonctionnement de l'agent \textbf{Mm.\
Gabriel}, agent multiloop local dédié et spécialiste de la géométrie du spectre
des nombres premiers, tel qu'analysé dans le dépôt local
\texttt{agent-multiloop-Gabriel-local} (version v3.35, licence Apache 2.0).

\subsection{Vocation et portée}

Gabriel est un \textbf{assistant mathématique expert HOL/Lean} entièrement
dédié à un objectif unique : assister l'auteur, Philippe Thomas Savard, à
construire l'outil géométrique dynamique --- la géométrie du spectre des
nombres premiers --- permettant d'apporter une réponse à l'énigme de Bernhard
Riemann et à la conjecture de la fonction zêta de Riemann. L'agent
\textbf{ne répond qu'aux questions portant sur la géométrie du spectre des
nombres premiers} ; toute question hors de ce champ est déclinée poliment. Sa
compétence principale est le fichier de preuve formelle
\texttt{theories/methode\_spectral.thy} (Isabelle/HOL ; 3\,714 lignes,
133 lemmes, 27 théorèmes), sur lequel toute réponse est ancrée.

\subsection{Architecture cognitive --- les 7 moteurs collaboratifs}

Le pipeline cognitif complet (\texttt{src/core/pipeline.py}) enchaîne sept
moteurs :
\[
\begin{gathered}
\texttt{RequestDecomposer}\to\texttt{ProofPlanner}\\
\to\texttt{SpectralCore}\to\texttt{RefinementLoop}\\
\to\texttt{CoherenceDetector}\to\texttt{SilentAuditLoop}
\end{gathered}
\]
où \texttt{RefinementLoop} assure l'auto-critique multiloop et
\texttt{SilentAuditLoop} l'audit silencieux anti-hallucination inversé. Un
\textbf{pré-raisonneur dynamique} (\texttt{PreReasoner}, fichier
\url{src/multiloop/pre_reasoner.py})
décide en T0, avant toute exécution, du niveau d'effort :

\begin{center}\small
\begin{tabular}{@{}l c l@{}}
\toprule
\rowcolor{spectralgris}
Mode & Itérations & Usage \\
\midrule
INSTANTANÉ & 0 & salutations, oui/non, continuation \\
RAPIDE & 1 & questions verbales/discursives sur Isabelle \\
STANDARD & 2 & calculs simples (RsP, gap, reconstruction) \\
APPROFONDI & 3 & configurations $n \times n$, blocs A/B \\
TRÈS COMPLEXE & 4 & théorie avancée, Pont Savard, Riemann \\
\bottomrule
\end{tabular}
\end{center}

Commandes CLI de forçage : \texttt{/rapide}, \texttt{/standard},
\texttt{/approfondi}, \texttt{/complet}, \texttt{/instantane}. Un mode
\og cinématique \fg{} affiche en temps réel mode détecté, itérations prévues,
chronomètre et ETA (rafraîchissement 2 Hz).

\subsection{Pipeline convolutif spectral (niveaux 1 et 2)}

Le dossier \texttt{pipeline\_cognitif/} implante le système convolutif décrit
aux sections précédentes --- c'est le moteur qui exécute l'arborescence
algorithmique de la section~5 :

\begin{center}\small
\begin{tabular}{@{}p{0.42\textwidth} p{0.54\textwidth}@{}}
\toprule
\rowcolor{spectralgris}
Module & Rôle \\
\midrule
\url{suites_geometriques_niveau2.py} & cœur : calculateurs niveaux 1 et 2, Digamma, formes fermées, ancrages XIV \\
\url{reconstruction_premiers_1_sur_k.py} & reconstructeur de référence 1/7 (16\,519 via $\SA - 7^{8}$) \\
\url{gabriel_geometric_wrapper_v74.py} & \texttt{PipelineCognitifNiveaux} : wrapper unifié (typique 1/2, non typiques $1/k$) \\
\url{Suites_geometriques_AB.py} & suites A/B géométriques générales \\
\url{fallback_geometrique.py} & repli géométrique si l'approche algébrique échoue \\
\url{metaphore_geometrique.py} & bloc \og Structure Géométrique Spatiale \fg{} des réponses \\
\url{validation_section_xiv.py} & validation de conformité à la Section XIV (exit 0 = conforme) \\
\url{multi_ratio_dispatcher.py} & distribution des requêtes multi-rapports $1/k$ \\
\bottomrule
\end{tabular}
\end{center}

\textbf{Niveau 1 (entier)} : suites A et B à termes entiers ($k^{i}$), règles
Savard terme à terme, 4 possibilités Digamma, division par $k^{6}$.
\textbf{Niveau 2 (géométrique)} : mêmes suites en forme $\sqrt{a^{2}+b^{2}}$
avec $a_1 = \sqrt{1+k^{2}}$ ; chaque terme géométrique vaut exactement
$a_1 \times$ terme entier, donc le facteur s'annule dans
$(\SB - \text{Digamma})/k^{6}$ : les deux niveaux ont les mêmes utilités.

\begin{boiteformule}[Conformité à la Section XIV de \texttt{methode\_spectral.thy} (0 écart)]
\begin{itemize}[nosep]
  \item Ancrages $n=10$ : rangs officiels $=$ rangs réels dans $\mathbb{P}$ (XIV.6) --- $k=2..9$ conformes (29, 227, 947, 2999, 7529, 16519, 32327, 58337) ;
  \item Formes fermées universelles $\SA$, $\SB$ (XIV.2) --- exactes pour $n \ge 8$ (arithmétique \texttt{Fraction}) ;
  \item Construction terme à terme (XIV.4) --- niveau 1 exact pour tout $n$ ; niveau 2 $= a_1 \times$ XIV.4 ;
  \item Validation exacte $k=8$, $n=34$ (XIV.7) --- $\SA$, $\SB$, Digamma conformes au chiffre près ; $P = 32\,537$ (rang 3\,492) ;
  \item Décalage de rang $\text{rang}(n) = \text{rang\_ancre} + (n-10)$ (XIV.6) --- $k=3$, $n=9..17$ : 223, 227, 229, 233, 239, 241, 251, 257, 263 ;
  \item Niveau 2 : mêmes ancrages que le niveau 1 (XIV.5) --- 8/8 conformes.
\end{itemize}
\end{boiteformule}

\subsection{Pipeline HOL du dépôt et corpus HOL}

Le pipeline HOL repose sur une \textbf{architecture concentrique à trois
sphères} formalisée dans \texttt{methode\_spectral.thy} et
\texttt{validation\_hol\_unifiee.thy} :
\[
\text{Ensemble} = \frac{1}{ms} + \frac{1}{t} + \frac{1}{x},
\qquad
\begin{cases}
1/x & \text{fonction zêta de Riemann (sphères } y_1, y_2, y_3\text{)},\\
1/t & \text{équation psi\_savard (Pont Chebyshev} \leftrightarrow \text{Spectral)},\\
1/ms & \text{Méthode Spectrale (sphères } ms_1, ms_2, ms_3\text{)}.
\end{cases}
\]

Trois concordances verrouillent la preuve : \textbf{C1} : $1/y_1 = 1/t$
(Chebyshev $=$ psi\_savard, validé numériquement) ; \textbf{C2} :
$1/y_3 = 1/ms_1$ (zéros $\leftrightarrow$ positions des premiers, exclusion des
composés) ; \textbf{C3} : $1/y_2 = 1/ms_3$ ($\mathrm{Re}(\rho) = 1/2 = R_{sP} = 1/2$,
alignement central). Théorème final : $R_{sP} = \mathrm{Re} = 1/2$, vrai sur tous
les entiers positifs distincts.

\begin{center}\small
\begin{tabular}{@{}p{0.42\textwidth} p{0.54\textwidth}@{}}
\toprule
\rowcolor{spectralgris}
Composante du corpus HOL & Caractéristiques \\
\midrule
\url{theories/methode_spectral.thy} & preuve principale : 3\,714 lignes, 133 lemmes, 27 théorèmes, 13 régimes, Section XIII / Pont Savard \\
\url{theories/validation_hol_unifiee.thy} & contre-validation jumelle HOL \\
\url{theories/validation_zeta_lean.thy} & contre-validation Lean/mathlib4 \\
\url{theories/ROOT} & session Isabelle \texttt{Methode\_Spectral} (Isabelle/HOL 2025-2) \\
\url{corpus/formulas_1_4.json}, \url{verif_p1097_n23_ratio4.thy} & corpus de formules et vérifications ciblées \\
\url{memory/dictionnaire_spectral.py} & 13 régimes cognitifs codifiés, dont \texttt{regime\_pont\_savard} \\
\url{memory/methode_spectral_section_XIII.py} & encyclopédie du Pont Savard \\
\url{gabriel_corpus.db} & base du corpus cognitif local (mémoire RAG) \\
\bottomrule
\end{tabular}
\end{center}

Une \textbf{factorisation formelle} par \emph{locale paramétré}
\texttt{spectral\_family} unifie les modèles 1/2, 1/3 et 1/4 ; trois
interprétations (\texttt{regime\_1\_2}, \texttt{regime\_1\_3},
\texttt{regime\_1\_4}) héritent du théorème universel
\texttt{RsP\_generic\_constant}, et l'extension à un modèle $1/k$ arbitraire ne
requiert qu'une seule ligne --- c'est le point d'ancrage HOL du système
convolutif multi-rapports du présent document. La section \og Foundations /
Meta-theory \fg{} du fichier \texttt{.thy} documente 6 postulats (P1..P6), les
3 opérations fondamentales (Reconstruction, Exclusion, Rapport spectral), la
règle Savard ($\text{Ensemble} = 1 = 1/x + 1/t + 1/ms$) et les 3 concordances
C1, C2, C3.

\subsection{Robustesse, audit et utilités pour la géométrie du spectre}

\begin{itemize}[nosep]
  \item \textbf{Tests} : 1\,702 tests Pytest (multiloop, RAG, syntaxe Isabelle, pipeline, PreReasoner, foundations, locale \texttt{spectral\_family}, Pont Savard, Section XIII) ; compilation Isabelle/HOL 2025-2 validée localement (Cygwin, 1\,min\,35) et en CI GitHub Actions (37 s, cache HIT) avec attestation SLSA (SHA256).
  \item \textbf{Audit} : \texttt{AuditStore}, trace JSON signée de chaque question/réponse ; chaque réponse est ancrée sur \texttt{methode\_spectral.thy} (empreinte SHA256 tracée).
  \item \textbf{Reconstruction} : calcul de $R_{sP}$, de \texttt{psi\_savard}, reconstruction du $n$-ième premier, exclusion des composés, en arithmétique entière exacte (\texttt{Fraction}/\texttt{sympy}) --- indispensable dès $k \ge 50$ où $\SA, \SB > 10^{15}$.
  \item \textbf{Contrat HOL} : le moteur produit les candidats C ; la validation HOL générale décide \texttt{EXCLU\_HOL}, \texttt{ANCRAGE\_POSSIBLE}, \texttt{P\_CERTIFIÉ} ou \texttt{C\_NON\_DÉCIDÉ} ; seul \texttt{P\_CERTIFIÉ} peut être annoncé comme nombre premier reconstruit (aucune sélection silencieuse).
  \item \textbf{Utilité spectrale} : le pipeline convolutif fournit à la géométrie du spectre un mécanisme constructif de localisation des premiers (ancrages $n=10$ + loi de rang spectral), tandis que le corpus HOL garantit la contre-validation formelle de chaque résultat annoncé.
\end{itemize}

\clearpage
\section*{Transcription corrig\'ee du PDF source : pages 22 \`a 48}
Les sections ci-dessous reprennent le texte des pages 22 a 48 du document source. La
transcription conserve le contenu et l'ordre du texte; seules les erreurs de lecture
\'evidentes et l'orthographe ont \'et\'e normalis\'ees. Les formules sont reproduites en
texte pour garder la transcription compl\`ete et comparable \`a la source.
\smallskip
\subsubsection*{Page 22}
\begin{Verbatim}[fontsize=\footnotesize]
2353
49  x  7-n - 7
6 = Somme suite A négative. 
*Quand n est un entier et qu'il est strictement négatif. 
 
Déterminer l'équation pour la somme de la suite B : 
(
2353
7
x  x  710) - Reste = 2260645142 Somme suite B. 
2353
7
2260645142
710
= 42.00000078  =>  x = 42 
(
2353
7
42  x  710) - Reste = 2260762792.167 
2260762792.167 - 2260645142 = 117650.167 = 705901
6 = Reste 
705901
6
7
6
= 100843 
Équation B n = 10 : 
(
2353
7
42  x  710) - 100843 (7
6) = 2260645142 Somme suite B. 
Forme générale pour la suite B (1/7) : 
(2353
294  x  7n) - 100843 (7
6) = Somme suite B. 
3.  SYSTÈME D'ÉQUATIONS CONVOLUTIF 
GÉNÉRALISÉ 
Pour tout rapport 1/k (k >= 2) et tout n commun 
3.1. NIVEAU 1  --  Constantes Savard universelles 
alphaA(k) = 2(k4 - k2 + 1) / ((k - 1)  .  k3) 
alphaB(k) = k  .  alphaA(k) 
deltaA(k) = k / (k - 1) 
deltaB(k) = (k7 - k6 + k) / (k - 1) 
\end{Verbatim}
\medskip
\subsubsection*{Page 23}
\begin{Verbatim}[fontsize=\footnotesize]
3.2. NIVEAU 2  --  Formules fermées (n quelconque, n commun à 
tous les rapports) 
SA(k, n) = (alphaA(k) / 2)  .  kn - deltaA(k) 
SB(k, n) = (alphaB(k) / 2)  .  kn - deltaB(k) 
Invariant spectral prouvé algébriquement (Théorème XIV .3) : 
(SA(k,n1) - SA(k,n2)) / (SB(k,n1) - SB(k,n2)) = 1/k for all n1 != n2, k >= 2 
3.3. NIVEAU 3  --  Construction terme à terme (Suite A, n >= 8) 
Position i Terme A(k,n,i) Terme B(k,n,i) 
1 <= i <= n-2 ki ki (i<=5) ; k7 (i=6, Saut Zêta) ; k^(i+1) (7<=i<=n-2) 
i = n-1 k^(n-1) - k^(n-3) k^n - k^(n-2) 
i = n k^n - k^(n-2) k^(n+1) - k^(n-1) 
 
3.4. NIVEAU 4  --  Mécanisme d'ancrage à n=10 (4 candidats 
Digamma) 
Les quatre essais dans l'ordre déterministe : 
# Position Signe Digamma Candidat P 
1 8 - SA(k,10) - k8 (SB - Digamma) / k6 
2 8 + SA(k,10) + k8 (SB - Digamma) / k6 
3 7 + SA(k,10) + k7 (SB - Digamma) / k6 
4 7 - SA(k,10) - k7 (SB - Digamma) / k6 
 
 ->  P certifié = premier candidat entier premier dans cet ordre. 
P(1/k, n=10) = (SB(k,10) - [SA(k,10) ± k7 ou ± k8]) / k6 
3.5. NIVEAU 5  --  Reconstruction pour n != 10 (n commun) 
rang_cible(k,n) = rang_base(k) + n - 10 
P(k,n) = premier rang_cible 
Digamma(k,n) = SB(k,n) - P(k,n)  x  k6 
\end{Verbatim}
\medskip
\subsubsection*{Page 24}
\begin{Verbatim}[fontsize=\footnotesize]
 
3.6. [OK]  VÉRIFICATION k=3 à 9 (tous corrects, 100%) 
k P(n=10) HOL Branche Calculé [OK] 
3 227 pos8 - 227    
4 947 pos8 + 947    
5 2 999 pos7 + 2 999    
6 7 529 pos8 + 7 529    
7 16 519 pos8 - 16 519    
8 32 327 pos8 - 32 327    
9 58 337 pos7 - 58 337    
 
3.7. TABLEAU COMPLET P(1/k), n=10  --  k=10 à k=111 
Légende :  P certifié unique |  C_MULTIPLE (plusieurs premiers, premier en ordre retenu) |  
AUCUN P (règle spéciale requise) 
k P(1/k, n=10) État Branche 
10 99 109     C_MULTIPLE pos8 - 
11  --    RÈGLE SPÉCIALE (PDF: 1 611 851)  --  
12 247 259    pos8 - 
13 368 939    pos8 + 
14  --    AUCUN P  --  
15 755 789    pos8 + 
16 1 044 751     C_MULTIPLE pos8 - 
17 1 414 943    pos7 + 
18 1 883 429    pos8 + 
19 2 469 277    pos7 - 
20 3 192 419     C_MULTIPLE pos8 - 
21 4 074 419     C_MULTIPLE pos8 + 
22  --    AUCUN P  --  
23 6 424 727    pos8 - 
\end{Verbatim}
\medskip
\subsubsection*{Page 25}
\begin{Verbatim}[fontsize=\footnotesize]
24 7 949 399     C_MULTIPLE pos8 - 
25 9 750 649     C_MULTIPLE pos8 - 
26 11 863 799    pos7 + 
27  --    RÈGLE SPÉCIALE HOL (P=14 330 707)  --  
28  --    AUCUN P  --  
29  --    AUCUN P  --  
30 24 273 929    pos8 - 
31  --    AUCUN P  --  
32 33 522 719    pos8 - 
33  --    AUCUN P  --  
34  --    AUCUN P  --  
35  --    AUCUN P  --  
36 60 420 851     C_MULTIPLE pos8 - 
37 69 293 303    pos7 + 
38  --    AUCUN P  --  
39  --    AUCUN P  --  
40 102 337 639     C_MULTIPLE pos8 - 
41  --    AUCUN P  --  
42 130 617 143    pos7 + 
43 146 929 021    pos7 - 
44  --    AUCUN P  --  
45  --    AUCUN P  --  
46 205 867 801     C_MULTIPLE pos8 - 
47  --    AUCUN P  --  
48  --    AUCUN P  --  
49 282 357 599     C_MULTIPLE pos7 + 
50  --    AUCUN P  --  
51  --    AUCUN P  --  
52 380 066 179    pos8 - 
53  --    AUCUN P  --  
\end{Verbatim}
\medskip
\subsubsection*{Page 26}
\begin{Verbatim}[fontsize=\footnotesize]
54 459 010 529     C_MULTIPLE pos8 - 
55  --    AUCUN P  --  
56  --    AUCUN P  --  
57 601 506 977    pos7 - 
58 656 165 077    pos8 - 
59  --    AUCUN P  --  
60  --    AUCUN P  --  
61  --    AUCUN P  --  
62 915 894 503    pos7 + 
63 992 182 589     C_MULTIPLE pos8 + 
64 1 073 483 839     C_MULTIPLE (3 premiers!) pos8 - 
65 1 160 020 289     C_MULTIPLE pos8 - 
66 1 252 049 501     C_MULTIPLE pos8 - 
67  --    AUCUN P  --  
68  --    AUCUN P  --  
69 1 563 702 839    pos7 + 
70 1 680 356 999    pos7 + 
71 1 803 876 551    pos8 - 
72  --    AUCUN P  --  
73  --    AUCUN P  --  
74  --    AUCUN P  --  
75 2 372 624 999    pos7 + 
76  --    AUCUN P  --  
77 2 706 333 629    pos8 - 
78 2 886 705 977    pos8 - 
79  --    AUCUN P  --  
80 3 276 294 479    pos8 - 
81  --    AUCUN P  --  
82  --    AUCUN P  --  
83  --    AUCUN P  --  
\end{Verbatim}
\medskip
\subsubsection*{Page 27}
\begin{Verbatim}[fontsize=\footnotesize]
84  --    AUCUN P  --  
85 4 436 438 999    pos7 + 
86  --    AUCUN P  --  
87 4 983 550 877    pos7 - 
88 5 276 645 527     C_MULTIPLE pos8 - 
89  --    AUCUN P  --  
90  --    AUCUN P  --  
91  --    AUCUN P  --  
92  --    AUCUN P  --  
93  --    AUCUN P  --  
94  --    AUCUN P  --  
95  --    AUCUN P  --  
96 8 152 842 431    pos7 - 
97 8 586 418 271     C_MULTIPLE pos8 + 
98  --    AUCUN P  --  
99  --    AUCUN P  --  
100  --    AUCUN P  --  
101  --    AUCUN P  --  
102 11 039 747 027    pos7 - 
103 11 591 637 509    pos8 + 
104  --    AUCUN P  --  
105 12 761 657 999    pos7 + 
106 13 381 076 101    pos8 - 
107 14 024 303 819    pos8 - 
108 14 692 021 271    pos7 - 
109 15 384 932 747    pos8 + 
110  --    AUCUN P  --  
111 16 849 226 351     C_MULTIPLE pos8 - 
 
\end{Verbatim}
\medskip
\subsubsection*{Page 28}
\begin{Verbatim}[fontsize=\footnotesize]
3.8.  Synthèse statistique (k=10 à 111) 
État Nombre de k % 
   P certifié unique 38 37 % 
    C_MULTIPLE (>=2 premiers) 20 20 % 
  AUCUN P (règle standard) 44 43 % 
 
3.9.  Notes importantes pour la programmation de l'agent 
Cas AUCUN P (44 valeurs)  --  Les 4 candidats standards sont tous composites. Deux 
stratégies possibles selon votre HOL : 
Règle spéciale cataloguée (comme k=27 : Digamma = SA - (2k8-k6), P=14 330 707) 
Mode CRÉATION D'ANCRAGE : l'agent retourne C_NON_CERTIFIÉ et attend validation HOL 
externe 
Cas C_MULTIPLE (20 valeurs)  --  Selon le contrat HOL : retourner tous les C/P possibles, 
aucune sélection silencieuse. L'agent doit les lister tous et attendre confirmation. 
k=11 est un cas forme spéciale : per votre fichier Moteur Géo, P=1 611 851 (rang 121 982) 
avec règle A7- documentée dans le PDF  --  non reconstruit par les 4 essais standards à 
n=10. 
Arithmétique exacte requise : pour k >= 50, les valeurs de SA et SB dépassent 10¹5  --  l'agent 
doit utiliser l'arithmétique entière Python (pas float) pour éviter les erreurs d'arrondi Excel 
4.  EXEMPLE NUMÉRIQUE COMPLET  --  n = 15 
4.1.  RAPPORT 1/13 (k = 13, n = 15) 
4.1.1. NIVEAU 1  --  Constantes Savard universelles 
Constante Formule Valeur exacte 
alpha_A(13) 2(134-13²+1) / (12·13³) 28393/13182 
alpha_B(13) 13·alpha_A(13) 28393/1014 
delta_A(13) 13/12 13/12 
delta_B(13) (137-136+13)/12 57 921 721/12 
 
\end{Verbatim}
\medskip
\subsubsection*{Page 29}
\begin{Verbatim}[fontsize=\footnotesize]
 
4.1.2. NIVEAU 2  --  Formules fermées 
Code 
Copier 
SA(13, 15) = (k4-k²+1)·k¹² / (k-1) - k/(k-1) 
           = 28 393  x  23 298 085 122 481 / 12 - 13/12 
           = 55 125 210 906 883 585 
SB(13, 15) = (k4-k²+1)·k¹³ / (k-1) - (k7-k6+k)/(k-1) 
           = 716 627 741 784 659 809 
Invariant spectral (Théorème XIV .3) : 
Code 
Copier 
DeltaSA = SA(n=15) - SA(n=10) = 55 125 210 906 883 585 - 148 468 220 264 
    = 55 125 062 438 663 321 
DeltaSB = SB(n=15) - SB(n=10) = 716 627 741 784 659 809 - 1 930 082 036 636 
    = 716 625 811 702 623 173 
DeltaSA / DeltaSB = 1/13   
4.1.3. NIVEAU 3  --  Construction terme à terme (n = 15) 
i A(13, 15, i) B(13, 15, i) Règle A Règle B 
1 13 13 k¹ k¹ 
2 169 169 k² k² 
3 2 197 2 197 k³ k³ 
4 28 561 28 561 k4 k4 
5 371 293 371 293 k5 k5 
6 4 826 809 62 748 517 k6 k7  <-  Saut Zêta 
7 62 748 517 815 730 721 k7 k8 
8 815 730 721 10 604 499 373 k8 k9 
9 10 604 499 373 137 858 491 849 k9 k¹0 
10 137 858 491 849 1 792 160 394 037 k¹0 k¹¹ 
11 1 792 160 394 037 23 298 085 122 481 k¹¹ k¹² 
12 23 298 085 122 481 302 875 106 592 253 k¹² k¹³ 
\end{Verbatim}
\medskip
\subsubsection*{Page 30}
\begin{Verbatim}[fontsize=\footnotesize]
13 302 875 106 592 253 3 937 376 385 699 289 k¹³ k¹4 
14 3 914 078 300 576 808 50 883 017 907 498 504 k¹4-k¹² k¹5-k¹³ 
15 50 883 017 907 498 504 661 479 232 797 480 552 k¹5-k¹³ k¹6-k¹4 
 
Code 
Copier 
Sigma A = 55 125 210 906 883 585   = SA(13,15) 
Sigma B = 716 627 741 784 659 809  = SB(13,15) 
4.1.4. NIVEAU 4  --  4 essais Digamma (k^6 = 4 826 809) 
 Les 4 essais sont calibrés pour n=10. Pour n=15, ils servent de vérification structurelle  --  
tous composites est le comportement attendu. 
Branche Digamma P candidat État 
pos13- 54 822 335 800 291 332 137 110 336 453   composé 
pos13+ 55 428 086 013 475 838 136 984 839 419   composé 
pos12+ 55 148 508 992 006 066 137 042 761 127   composé 
pos12- 55 101 912 821 761 104 137 052 414 745   composé 
 
4.1.5. NIVEAU 5  --  Reconstruction rang cible (méthode officielle pour n 
!= 10) 
Code 
Copier 
P(13, n=10) documenté = 368 939 
rang_base(13)  = primepi(368 939) = 31 452 
Deltan = 15 - 10 = 5 
rang_cible(13, n=15) = 31 452 + 5 = 31 457 
P(13, n=15) = prime(31 457) = 369 023   PREMIER certifié 
Digamma (13,15) = SB - P  x  k6 
               = 716 627 741 784 659 809 - 369 023  x  4 826 809 
               = 716 625 960 581 122 202 
\end{Verbatim}
\medskip
\subsubsection*{Page 31}
\begin{Verbatim}[fontsize=\footnotesize]
4.2.  RAPPORT 1/42 (k = 42, n = 15) 
4.2.1. NIVEAU 1  --  Constantes Savard universelles 
Constante Valeur exacte 
alpha_A(42) 3 109 933 / 1 518 804 
alpha_B(42) 3 109 933 / 36 162 
delta_A(42) 42/41 
delta_B(42) 225 050 301 546 / 41 
 
4.2.2. NIVEAU 2  --  Formules fermées 
Code 
Copier 
SA(42, 15) = 2 285 381 254 365 702 554 104 806 
SB(42, 15) = 95 986 012 683 359 501 783 370 150 
DeltaSA / DeltaSB = 1/42   
4.2.3. NIVEAU 3  --  Construction terme à terme (n = 15) 
i A(42, 15, i) B(42, 15, i) 
1 42 42 
2 1 764 1 764 
3 74 088 74 088 
4 3 111 696 3 111 696 
5 130 691 232 130 691 232 
6 5 489 031 744 230 539 333 248  <-  Saut Zêta (k7) 
7 230 539 333 248 9 682 651 996 416 
8 9 682 651 996 416 406 671 383 849 472 
9 406 671 383 849 472 17 080 198 121 677 824 
10 17 080 198 121 677 824 717 368 321 110 468 608 
11 717 368 321 110 468 608 30 129 469 486 639 681 536 
12 30 129 469 486 639 681 536 1 265 437 718 438 866 624 512 
13 1 265 437 718 438 866 624 512 53 148 384 174 432 398 229 504 
14 53 118 254 704 945 758 547 968 2 230 966 697 607 721 859 014 656 
\end{Verbatim}
\medskip
\subsubsection*{Page 32}
\begin{Verbatim}[fontsize=\footnotesize]
15 2 230 966 697 607 721 859 014 656 93 700 601 299 524 318 078 615 552 
 
Code 
Copier 
Sigma A = 2 285 381 254 365 702 554 104 806   
Sigma B = 95 986 012 683 359 501 783 370 150   
4.2.4. NIVEAU 4  --  4 essais Digamma (k^6 = 5 489 031 744) 
Branche Digamma P candidat État 
pos13- 2 284 115 816 647 263 687 480 294 17 070 751 498 046 399   composé 
pos13+ 2 286 646 692 084 141 420 729 318 17 070 290 419 379 903   composé 
pos12+ 2 285 411 383 835 189 193 786 342 17 070 515 469 681 407   composé 
pos12- 2 285 351 124 896 215 914 423 270 17 070 526 447 744 895   composé 
 
4.2.5. NIVEAU 5  --  Reconstruction rang_cible 
Code 
Copier 
P(42, n=10) documenté = 130 617 143 
rang_base(42)  = primepi(130 617 143) = 7 411 117 
Deltan = 15 - 10 = 5 
rang_cible(42, n=15) = 7 411 117 + 5 = 7 411 122 
P(42, n=15) = prime(7 411 122) = 130 617 251   PREMIER certifié 
Digamma (42,15) = 95 986 012 683 359 501 783 370 150 
               - 130 617 251  x  5 489 031 744 
               = 95 986 011 966 397 264 730 354 406 
4.3.  RAPPORT 1/54 (k = 54, n = 15) 
4.3.1. NIVEAU 1  --  Constantes Savard universelles 
Constante Valeur exacte 
alpha_A(54) 8 500 141 / 4 172 796 
alpha_B(54) 8 500 141 / 77 274 
delta_A(54) 54/53 
delta_B(54) 1 314 130 298 742 / 53 
\end{Verbatim}
\medskip
\subsubsection*{Page 33}
\begin{Verbatim}[fontsize=\footnotesize]
 
4.3.2. NIVEAU 2  --  Formules fermées 
Code 
Copier 
SA(54, 15) = 98 599 651 085 954 948 762 648 034 
SB(54, 15) = 5 324 381 158 641 567 208 388 082 594 
DeltaSA / DeltaSB = 1/54   
4.3.3. NIVEAU 3  --  Construction terme à terme (n = 15) 
i A(54, 15, i) B(54, 15, i) 
1-
5 
54 / 2 916 / 157 464 / 8 503 056 / 459 165 
024 
idem 
6 24 794 911 296 1 338 925 209 984  <-  Saut Zêta (547) 
7 1 338 925 209 984 72 301 961 339 136 
8 72 301 961 339 136 3 904 305 912 313 344 
9 3 904 305 912 313 344 210 832 519 264 920 576 
10 210 832 519 264 920 576 11 384 956 040 305 711 104 
11 11 384 956 040 305 711 104 614 787 626 176 508 399 616 
12 614 787 626 176 508 399 616 33 198 531 813 531 453 579 264 
13 33 198 531 813 531 453 579 264 1 792 720 717 930 698 493 280 256 
14 1 792 105 930 304 521 984 880 640 96 773 720 236 444 187 183 554 560 
15 96 773 720 236 444 187 183 554 560 5 225 780 892 767 986 107 911 946 
240 
 
Code 
Copier 
Sigma A = 98 599 651 085 954 948 762 648 034   
Sigma B = 5 324 381 158 641 567 208 388 082 594   
\end{Verbatim}
\medskip
\subsubsection*{Page 34}
\begin{Verbatim}[fontsize=\footnotesize]
 
 
4.3.4. NIVEAU 4  --  4 essais Digamma (k^6 = 24 794 911 296) 
Branche Digamma P candidat État 
pos13- 98 566 452 554 141 417 309 068 
770 
210 761 581 023 702 
719 
  
composé 
pos13+ 98 632 849 617 768 480 216 227 
298 
210 758 903 173 282 
751 
  
composé 
pos12+ 98 600 265 873 581 125 271 047 
650 
210 760 217 303 581 
439 
  
composé 
pos12- 98 599 036 298 328 772 254 248 
418 
210 760 266 893 404 
031 
  
composé 
 
4.3.5. NIVEAU 5  --  Reconstruction rang_cible 
Code 
Copier 
P(54, n=10) documenté = 459 010 529 
rang_base(54)  = primepi(459 010 529) = 24 304 918 
Deltan = 15 - 10 = 5 
rang_cible(54, n=15) = 24 304 918 + 5 = 24 304 923 
P(54, n=15) = prime(24 304 923) = 459 010 597   PREMIER certifié 
Digamma(54,15) = 5 324 381 158 641 567 208 388 082 594 
               - 459 010 597  x  24 794 911 296 
               = 5 324 381 147 260 440 171 849 078 882 
4.4.  TABLEAU RÉCAPITULATIF FINAL  --  P(1/k, n = 15) 
k SA(k, 15) SB(k, 15) rang_base rang_cible P(k, 
n=15) 
P(k, 
n=10) 
État 
13 55 125 210 
906 883 
585 
716 627 741 
784 659 809 
31 452 31 457 369 
023 
368 
939 
   
CERTIFIÉ 
42 2 285 381 
254 365 
95 986 012 
683 359 501 
783 370 150 
7 411 117 7 411 122 130 
617 
251 
130 
617 
143 
   
CERTIFIÉ 
\end{Verbatim}
\medskip
\subsubsection*{Page 35}
\begin{Verbatim}[fontsize=\footnotesize]
702 554 
104 806 
54 98 599 651 
085 954 
948 762 
648 034 
5 324 381 
158 641 567 
208 388 082 
594 
24 304 
918 
24 304 923 459 
010 
597 
459 
010 
529 
   
CERTIFIÉ 
 
4.5.  Notes critiques pour l'agent 
4.5.1. 1. Saut Zêta  --  asymétrie A / B confirmée 
La correction importante validée ici : A(k, n, i) = k^i pour tout 1 <= i <= n-2 (pas de Saut
    Zêta
dans la suite A). Le Saut Zêta k7 à la position i=6 n'existe que dans la suite B. Les deux
    suites
sont identiques pour i <= 5, puis divergent à partir de i = 6. 
4.5.2. 2. Niveau 4 hors-calibrage pour n != 10 
Les 4 essais Digamma (pos(n-2)±, pos(n-3)±) produisent tous des candidats composites 
pour n=15  --  ce comportement est attendu et correct. Ces essais sont calibrés 
structurellement sur n=10 ; le Niveau 5 est la voie officielle pour tout n != 10. 
4.5.3. 3. Niveau 5  --  formule rang_cible 
python 
Copier 
rang_base(k)       = primepi(P(k, n=10))   # rang du P certifié à n=10 
rang_cible(k, n)   = rang_base(k) + (n - 10) 
P(k, n)            = prime(rang_cible(k, n))  # arithmétique exacte sympy 
L'écart en rang +5 (pour Deltan = 5) est cohérent pour les trois k, confirmant l'universalité de
    la
loi de rang spectral. 
4.5.4. 4. Invariant spectral (Théorème XIV .3)  --  vérifié pour les 3 rapports 
Code 
Copier 
DeltaSA(k, n=15, n=10) / DeltaSB(k, n=15, n=10) = 1/k   for all k  in  {13, 42, 54}   
4.5.5. 5. Arithmétique  --  seuils de dépassement float 
k SA(k,15) Ordre de grandeur Float 64 safe= 
13 55 125 210 906 883 585 ~5.5 x 10¹6     limite 
42 2.29 x 10²4 ~10²4   erreur garantie 
54 9.86 x 10²5 ~10²5   erreur garantie 
\end{Verbatim}
\medskip
\clearpage
\subsubsection*{Page 36}
\begin{Verbatim}[fontsize=\footnotesize]
 
 ->  Arithmétique entière Python obligatoire dès k >= 42 pour n = 15. Excel / float64 
introduiraient des erreurs d'arrondi fatales. 
 
Comparaison structurelle entre les ancrages premiers de la géométrie du spectre des 
nombres premiers et les zéros non triviaux de la fonction zêta de Riemann 
1. Deux géométries distinctes des nombres premiers 
La fonction zêta de Riemann encode les nombres premiers dans le plan complexe à travers 
la relation fondamentale : 
zeta(s) = sum n-s
infinity
n=1
. 
Les zéros non triviaux de zeta(s)se situent dans la bande critique 0 < Re(s) < 1, et leur 
distribution est décrite par la formule asymptotique : 
N(T) ~ T
2pi log ( T
2pi) - T
2pi, 
où N(T)désigne le nombre de zéros non triviaux dont la partie imaginaire satisfait | t |<= T. 
Cette loi est presque linéaire en T, modulée par un facteur logarithmique. 
Dans la géométrie du spectre des nombres premiers, la situation est fondamentalement 
différente : il ne s'agit pas de compter des zéros, mais de reconstruire des nombres 
premiers à partir de rapports typiques et non typiques 1/k. Chaque rapport 1/kproduit un 
entier P(k), qui peut être un nombre premier ou un composé. 
Les rapports typiques et non typiques qui possèdent un ancrage pour n = 10 --  comme 
démontré précédemment  --  retournent un élément de l'ensemble complet des nombres 
premiers P. Inversement, lorsque le système convolutif (et l'équation qui permet la 
reconstruction des premiers pour des rapports multiples) ne trouve aucun ancrage, aucun 
nombre premier n'est retourné pour l'ensemble P. 
Le rapport typique 1/k = 1/2 retourne pour n = 10 l'ancrage P = 29. Dans ce cas, la valeur
de ncorrespond exactement à la position du nombre premier dans l'ordre chronologique 
des premiers. 
Pour les rapports non typiques 1/k != 1/2, la valeur de nreprésente uniquement la quantité 
de termes dans les suites Aet B. Le nombre premier retourné ne correspond pas à la 
position n, contrairement au cas typique. 
Ainsi : 
-  Rapport typique : 
n = quantite

 de termes dans A et B = position du premier. 
-  Rapport non typique : 
\end{Verbatim}
\medskip
\subsubsection*{Page 37}
\begin{Verbatim}[fontsize=\footnotesize]
n = quantite

 de termes dans A et B != position du premier. 
Il est important de rappeler que la position des nombres premiers dans la géométrie du 
spectre suit l'ordre chronologique croissant : 
2 = 1er,  3 = 2ie

me,  5 = 3ie

me,  7 = 4ie

me,  Pn = nie

me premier. 
La même convention s'applique aux nombres premiers négatifs : 
-2 = -1er,   - 3 = -2ie

me,   - 5 = -3ie

me,   - 7 = -4ie

me,   - Pn = -nie

me premier. 
Validation HOL/Isabelle et le "pont Savard" 
La géométrie du spectre des nombres premiers s'appuie sur une validation formelle 
réalisée à l'aide de l'assistant de preuve Isabelle et du code HOL contenu dans le fichier 
methode_spectral.thy. La section 13, intitulée « pont Savard », établit un lien entre la 
géométrie du spectre et l'énigme de Bernhard Riemann, en particulier la conjecture portant 
sur la fonction zêta. 
Cette section démontre, via un théorème local au fichier thy, que : 
RsP = Re(s) = 1
2 
est vrai dans le cadre de la validation methode_spectral.thy. 
Le « pont Savard » met en relation, à l'aide d'une équation, certains invariants communs 
entre : 
-  la géométrie du spectre des nombres premiers (validée dans le script HOL), 
-  et la fonction zêta de Riemann. 
L'équation de Tchebyshev joue un rôle central dans cette correspondance. Philippe Thomas 
Savard identifie cette équation comme un invariant fondamental partagé par les deux 
structures, en raison de son unicité au sein de la théorie analytique des nombres. 
Savard adapte l'équation de Tchebyshev à la géométrie du spectre, rendant les deux 
formulations difficilement distinguables. 
Voici la correspondance : 
Équation de Tchebyshev : 
psi(x) = x - xp
p - log(2pi) - 1
2 log(1 - x-2) 
Équation psi (Savard) : 
psiSavard(x) = x - kn
Sbn
- log(2pi) - 1
2 log(1 - x-2) 
Dans l'équation de Tchebyshev, le terme 
xp
p correspond, dans l'équation de Savard, à : 
\end{Verbatim}
\medskip
\subsubsection*{Page 38}
\begin{Verbatim}[fontsize=\footnotesize]
kn
Sbn
. 
où : 
-  kest la valeur du rapport typique 1/k = 1/2ou non typique 1/k != 1/2, 
-  Sbnest la somme de la suite B. 
Par exemple, pour le rapport typique 1/k = 1/2et n = 10, l'ancrage retourné est P = 29, et 
la somme de la suite Best : 
Sbn = 21 + 22 + 23 + 24 + 25 + 26 + 27 + 28 + (29 - 27) + (210 - 28) = 3262. 
 
La formulation de Tchebyshev adaptée par Savard commence par l'évaluation suivante : 
44
4 + 33
3 + 52
5 + 71
1 + 111
1 + 131
1 + 171
1 + 191
1 + 231
1 + 291
1 = 197, 
soit un total de 10 termes. 
On calcule ensuite : 
log(197)
log(10) = 2.294466226. 
Puis : 
2.294466226 - log(2pi) - 1
2 log(1 - 30-2) = 1.496527767. 
Enfin : 
30 - 1.496527767 = 28.50347223. 
La valeur généralement obtenue pour l'ancrage du premier 29 à l'aide de l'équation de
Tchebyshev est : 
28.49765. 
Comparaison avec l'équation psi (Savard) 
L'équation psi (Savard) donne pour le même cas : 
psiSavard = 30 - 210
3262 - log(pi) - 1
2 log(1 - 30-2) = 28.8881437. 
La valeur exacte recherchée étant 29, l'équation psi (Savard) présente une oscillation plus 
faible et une ondulation moins prononcée que celle de Tchebyshev. De plus, la précision de 
psiSavard augmente systématiquement lorsque la valeur du nombre premier considéré croît.
Par exemple, pour le nombre premier 227, qui est le 49e premier : 
psiSavard(227) = 226.894132. 
\end{Verbatim}
\medskip
\subsubsection*{Page 39}
\begin{Verbatim}[fontsize=\footnotesize]
La partie réelle 0.894132est plus proche de l'entier exact 227que ne l'est la partie réelle 
0.8881437obtenue pour 29. Elle est également plus précise que l'ondulation obtenue pour 
le premier 31, où : 
psiSavard(31) = 30.89125839. 
Cette amélioration constante de la précision montre que, à grande hauteur sur la droite 
critique, l'équation psi (Savard) devient asymptotiquement exacte, ce qui confirme la 
pertinence de son comportement pour des valeurs très grandes comme pour des valeurs 
proches de zéro. 
Lien avec la validation HOL/Isabelle 
L'équation psi (Savard), conçue comme une méta-analyse de l'équation de Tchebyshev, est 
mise en avant dans la section 13 du fichier methode_spectral.thy de la validation 
HOL/Isabelle. Cette section traite également de la répartition des nombres premiers et de 
leur correspondance avec les zéros de la droite critique de la fonction zêta. 
La répartition des premiers est au coeur de l'énigme de Riemann, qui pose la question : 
Re(s) = 1
2   = 
Le fichier methode_spectral.thy, conçu à l'aide de l'assistant de preuve Isabelle, inclut 
également une section intitulée « preuve par l'absurde », démontrant que les entiers 
composés Csont nécessairement exclus de la méthode spectrale. 
Théorème de la section "Preuve par l'absurde" (methode_spectral.thy, lignes 1982-
2002) 
isabelle 
text \<open> 
  Puisque prime_i i est défini via un choix de Hilbert sur la propriété
  "prime p \<and> position p = i" , et que prime_i_is_prime démontre que
  prime (prime_i i) tient toujours, il est logiquement impossible qu'un 
  entier composé C soit égal à prime_i i pour un i quelconque.
\<close> 
 
theorem composite_not_prime_i: 
  fixes C :: nat 
  assumes "~ prime C" 
  shows "ALL i. C ~= prime_i i" 
proof (rule allI, rule ccontr) 
\end{Verbatim}
\medskip
\subsubsection*{Page 40}
\begin{Verbatim}[fontsize=\footnotesize]
  fix i 
  assume "~ (C ~= prime_i i)" 
  hence eq: "C = prime_i i" by simp 
  have "prime (prime_i i)" by (rule prime_i_is_prime) 
  with eq have "prime C" by simp 
  with assms show False by contradiction 
qed 
Ce théorème démontre élégamment que si l'on attribue à un entier composé Cla qualité 
d'être un premier P, alors C = Pipour un certain i, ce qui est impossible. Ainsi, les 
composés sont exclus de la méthode spectrale de manière certaine. 
Cette exclusion se manifeste également dans l'historique de l'agent Gabriel multiloop, 
conçu par Savard à partir du corpus HOL et du pipeline cognitif incluant le système 
convolutif. Les logs de l'agent montrent que lorsqu'une requête demande un nombre 
composé, la méthode échoue, car elle ne s'applique qu'aux nombres premiers. 
Équation de reconstruction des premiers 
Somme Suite B - Digamma Calcule

t6(zeta) = Pi. 
Si l'on remplace Pipar un composé Ci, la valeur de nobtenue n'est plus un entier 
strictement positif (pour les premiers positifs) ni strictement négatif (pour les premiers 
négatifs). Elle contient une partie réelle, ce qui prouve que le résultat n'est pas un nombre 
premier. Ainsi, la méthode spectrale exclut automatiquement les composés. 
Portée de la méthode spectrale 
Pour les rapports typiques et non typiques 1/kavec k  in  [2, 10001], la méthode spectrale 
retourne 2604 ancrages premiers, surpassant largement les records actuels concernant 
les zéros non triviaux (~= 1012) en atteignant une échelle de : 
1020. 
Par exemple, pour le rapport non typique 1/9999 et n = 10, l'ancrage retourné est :
99950008999300019999. 
La méthode spectrale détermine donc 2604 rapports 1/k pour lesquels un ancrage premier 
est obtenu, chacun définissant un ensemble complet P. 
Correspondance avec la fonction zêta 
La fonction zêta, sa droite critique et ses zéros non triviaux trouvent un homologue 
structurel dans la géométrie du spectre des nombres premiers. Le rapport typique 1/k =
1/2représente la droite critique, tandis que les rapports non typiques correspondent aux 
hauteurs tronquées Tdans la formule : 
\end{Verbatim}
\medskip
\subsubsection*{Page 41}
\begin{Verbatim}[fontsize=\footnotesize]
N(T) ~ T
2pi log ( T
2pi) - T
2pi. 
Par exemple, le rapport non typique 1/3 retourne pour n = 10 le 49e nombre premier, soit
227. Dans ce cas, toutes les valeurs premières inférieures à 173sont considérées comme 
négatives pour ce rapport. 
De même, le rapport non typique 1/4 retourne pour n = 10 le premier 947, ce qui rend
toutes les valeurs premières inférieures à 881négatives pour ce rapport. Ainsi, un même 
nombre premier (par exemple 227) peut être positif pour un rapport non typique, et négatif 
pour un autre. 
Cette bipolarité (+/-)entre rapports non typiques constitue une caractéristique 
essentielle de la géométrie du spectre des nombres premiers. 
 
Les équations associées au rapport typique 1/k = 1/2, qui déterminent les sommes des
suites Aet Bet généralisent la valeur de npour les nombres premiers négatifs, sont les 
suivantes : 
Somme suite A = 3.25  x  2-n - 2, 
Somme suite B = 6.5  x  2-n - 66. 
L'équation psi (Savard), adaptée au cas négatif, s'écrit : 
psiSavard = -x - 2-n
Sbn
- log(2pi) - 1
2 log(1 - (-x)-2). 
Par exemple, pour x = -28et n = -10, avec Sbn = -65.99365234, on obtient : 
psiSavard = -28 - 2-10
-65.99365234 - log(2pi) - 1
2 log(1 - (-28)-2) = -28.79844187. 
La valeur exacte recherchée est -29, soit le -10e nombre premier. La valeur obtenue par 
psiSavard est plus proche de -29 que ne l'est la valeur positive correspondante obtenue pour
29, où : 
psiSavard(29) = 28.8881437. 
L'ondulation est donc plus précise pour le premier négatif -29. 
Écarts entre nombres premiers : positifs, négatifs et mixtes 
La géométrie du spectre des nombres premiers étudie également les écarts entre premiers. 
La validation HOL/Isabelle (methode_spectral.thy) distingue trois types d'écarts : 
1. Écarts positifs : entre deux premiers positifs. 
2. Écarts négatifs : entre deux premiers négatifs. 
3. Écarts mixtes : entre un premier positif et un premier négatif. 
Exemple : écart entre les premiers 23 et 7 
\end{Verbatim}
\medskip
\subsubsection*{Page 42}
\begin{Verbatim}[fontsize=\footnotesize]
Pour déterminer la quantité de nombres entre 23 et 7, on procède comme suit : 
1. Somme de la suite A du premier suivant 7, soit 11 
(7 est le 4e premier, 11 est le 5e.) 
Somme A(11) = (3.25
2  x  25) - 2 = 50. 
2. Somme de la suite B du premier 23 
Somme B(23) = (6.5
2  x  29) - 66 = 1598. 
3. Digamma calculé de 23 
(1598/64 - 23)  x  64 = 126. 
4. Calcul de l'écart 
50 - (1598 - 126) = -1442. 
5. Digamma calculé de 7 
Somme B(7) = (6.5
2  x  24) - 66 = -14, 
((-14)/64 - 7)  x  64 = -464. 
6. Quantité de nombres entre 23 et 7 
-1442 - (-462)
64 = -15. 
Les nombres entre 7 et 23 sont donc : 
8,9,10,11,12,13,14,15,16,17,18,19,20,21,22, 
soit 15 nombres. 
Exemple : écart entre les premiers négatifs -19 et -5 
1. Déterminer le premier précédent -19, soit -17 
(-17 est le 7e premier négatif.) 
Somme A(-17) = 3.25  x  2-7 - 2 = - 10110
5120 . 
2. Somme de la suite B du premier -5 
Somme B(-5) = 6.5  x  2-3 - 66 = - 20860
320 . 
3. Digamma calculé de -5 
(-20860/320
64 - (-5))  x  64 = 81540
320 . 
4. Calcul intermédiaire 
\end{Verbatim}
\medskip
\subsubsection*{Page 43}
\begin{Verbatim}[fontsize=\footnotesize]
- 10110
5120 - (- 20860
320 - 81540
320 ) = 1628290
5120 . 
5. Somme de la suite B de -19 
Somme B(-19) = 6.5  x  2-8 - 66 = - 337790
5120 . 
6. Digamma calculé de -19 
(-337790/5120
64 - (-11))  x  64 = 5888130
5120 . 
7. Quantité de nombres entre -19 et -5 
1628290/5120 - 5888130/5120
64 = -13. 
Les nombres entre -19 et -5 sont donc : 
-18, -17, -16, -15, -14, -13, -12, -11, -10, -9, -8, -7, -6, 
soit 13 nombres. 
Interprétation géométrique des écarts négatifs 
Pour les écarts positifs, la première étape consiste à déterminer le premier suivant la valeur 
minimale (par exemple, 7  ->  11). Pour les écarts négatifs, la première étape consiste à 
déterminer le premier précédent la valeur minimale (par exemple, -19  ->  -17). 
La géométrie du spectre adopte une convention idiomatique pour les premiers négatifs : 
-  -2est le plus grand des premiers négatifs, 
-  -3est le deuxième plus grand, 
-  puis -5, -7, etc. 
Cette convention est analogue à l'échelle des températures : 
-  -1 deg Cest plus chaud que -30 deg C, 
-  bien que les valeurs soient négatives. 
Ainsi, sur une droite représentant l'ensemble Ppositif et négatif : 
-  la valeur minimale est toujours située à gauche, 
-  la valeur suivante est à droite, 
-  que l'on soit dans le domaine positif ou négatif. 
Cette structure permet de traiter les écarts négatifs de manière cohérente avec les écarts 
positifs. 
 
\end{Verbatim}
\medskip
\subsubsection*{Page 44}
\begin{Verbatim}[fontsize=\footnotesize]
L 'équation psi (Savard) appliquée aux quinze premiers nombres premiers négatifs 
produit une ondulation dont la précision varie selon la valeur de Sbn. Au-delà du 
quinzième premier négatif, soit -47, la valeur de Sbndevient pratiquement égale à 
66. Ainsi, pour tous les nombres premiers négatifs précédant -47, l'ondulation 
demeure essentiellement identique, car les exposants contraignent la valeur de 
Sbnà se stabiliser autour de 66. 
Les quinze premiers nombres premiers négatifs  --  les plus grands de l'infini négatif 
 --  sont : 
-2, - 3, - 5, - 7, - 11, - 13, - 17, - 19, - 23, - 29, - 31, - 37, - 41, - 43, - 47. 
Ces quinze valeurs constituent la « tête » de l'infini premier négatif, tandis que toutes 
les valeurs précédentes sont plus petites. 
Limites numériques et nécessité d'une approche algébrique 
La méthode spectrale, lorsqu'elle utilise une approche numérique pour déterminer 
la valeur de n, présente une difficulté inhérente : l'exposant 2-nproduit des valeurs 
extrêmement grandes lorsque nest strictement positif, et des valeurs infinitésimales 
lorsque nest strictement négatif. 
Pour les très grands nombres premiers, l'approche numérique devient impraticable. 
Il faut alors recourir à une approche algébrique, ce que Savard juge 
conceptuellement délicat, car la question de Riemann porte précisément sur la 
nature analytique des zéros : 
« Est-ce que tous les zéros non triviaux de la fonction zêta de Bernhard Riemann ont 
pour partie réelle 1/2= » 
Pour Savard, une « partie réelle algébrique » ne peut être qu'une partie imaginaire, ce 
qui contredit l'esprit de la conjecture et ne peut donc servir à résoudre l'énigme. 
Suites et rapport spectral 
Les zéros de la fonction zêta sont définis par des suites, par exemple : 
delta(2) = pi2
6 . 
La méthode spectrale adopte une approche analogue : les suites Aet Bdéterminent 
la quantité de termes correspondant à la valeur de n, laquelle représente la position 
du nombre premier dans le cas du rapport typique. 
Les équations permettant de déterminer le rapport spectral RsP = 1/2 entre deux
premiers symétriques sont : 
Comparaison symétrique 1 x 1 
((3.25/2  x  2n1) - 2) - ((3.25/2  x  2n2) - 2)
((6.5/2  x  2n1) - 66) - ((6.5/2  x  2n2) - 66) = 1
2. 
\end{Verbatim}
\medskip
\subsubsection*{Page 45}
\begin{Verbatim}[fontsize=\footnotesize]
Cas incluant un premier négatif 
((3.25  x  2n1) - 2) - ((3.25/2  x  2n2) - 2)
((6.5  x  2n1) - 66) - ((6.5/2  x  2n2) - 66) = 1
2. 
Comparaison symétrique n x n 
((3.25/2  x  2n1) - 2) - ((3.25/2  x  2n2) - 2) - ((3.25/2  x  2n1) - 2) + ((3.25/2  x  2n2) -
    2)
((6.5/2  x  2n1) - 66) - ((6.5/2  x  2n2) - 66) - ((6.5/2  x  2n1) - 66) + ((6.5/2  x  2n2) -
    66)
= 1
2. 
Les équations des suites A et B permettent de calculer les sommes de manière
efficace, sans devoir dénombrer chaque terme individuellement. Elles sont exactes 
pour toute valeur entière de n. 
La coïncidence 3.25 / 6.5 n'est pas la raison du rapport spectral 
Considérons les équations : 
x + 2 = 13, 
2x + 4 = 109. 
On observe que : 
x + 2
2x + 109 != 1
2, 
bien que la simplification algébrique puisse donner l'illusion d'un rapport 1/2. De 
même, le fait que : 
3.25
6.5 = 1
2 
est une coïncidence algébrique, et non la raison profonde du rapport spectral 
RsP = 1/2. 
Exemple : premiers 7 et 31 
((3.25/2  x  24) - 2) - ((3.25/2  x  211) - 2)
((6.5/2  x  24) - 66) - ((6.5/2  x  211) - 66) = RsP. 
Ce qui donne : 
24 - 3326
-14 - 6590 = 1
2. 
Or : 
24
-14 != 1
2 ,  3326
6590 != 1
2. 
La simplification algébrique ne reflète donc pas la véritable structure du rapport 
spectral. 
\end{Verbatim}
\medskip
\clearpage
\subsubsection*{Page 46}
\begin{Verbatim}[fontsize=\footnotesize]
Évidence fondamentale de la géométrie du spectre 
« Lorsque n >= 1ou n <= -1et que nest un entier, toutes les valeurs de nramènent à 
un nombre premier P. Toutes les valeurs de nsont la conséquence directe de la 
quantité de termes dans les suites Aet B. Tous les nombres premiers Prespectent le 
rapport spectral 1/k. Ce rapport est numériquement valide, mais algébriquement 
incohérent. » 
Lien avec la fonction zêta et les rapports non typiques 
L 'axiome concernant les parties réelles algébriques introduit des difficultés dans la 
validation HOL/Isabelle. Cependant, grâce au pipeline cognitif et au corpus HOL 
intégré dans le dépôt GitHub, la géométrie du spectre permet de reconstruire les 
nombres premiers pour la valeur généralisée de n, pour tous les rapports typiques et 
non typiques. 
Le rapport typique 1/k = 1/2correspond à la droite critique de la fonction zêta. Les 
rapports non typiques sont l'analogue des zéros non triviaux. 
Dans la formule : 
N(T) ~ T
2pi log ( T
2pi) - T
2pi, 
la valeur Treprésente une hauteur tronquée de la droite critique. Les rapports non 
typiques 1/k != 1/2jouent un rôle similaire : la valeur de nne correspond pas à la 
position du nombre premier, mais uniquement à la quantité de termes dans les 
suites Aet B. 
Par exemple : 
-  Pour 1/k = 1/3 et n = 10, le premier retourné est 227, soit le 49e premier.
-  Pour 1/k = 1/4 et n = 10, le premier retourné est 947, ce qui rend toutes les
valeurs premières inférieures à 881 négatives pour ce rapport. 
Ainsi, un même nombre premier peut être positif pour un rapport non typique et 
négatif pour un autre, en raison de la polarité déterminée par la valeur de n. 
Les ancrages premiers observés dans la géométrie du spectre suivent une loi empirique de 
type puissance : 
log P (k) ~= a log k + b,  avec a ~= 5. 
Ce qui équivaut à : 
P(k) ~= C  .  k5. 
Il s'agit d'une loi polynomiale, et non logarithmique. Elle définit une géométrie alternative 
des nombres premiers, dans laquelle la taille des premiers reconstruits croît comme une 
puissance de l'indice k. 
2. Comparaison des ordres de grandeur 
\end{Verbatim}
\medskip
\subsubsection*{Page 47}
\begin{Verbatim}[fontsize=\footnotesize]
Les zéros non triviaux de la fonction zeta(s)atteignent des hauteurs de l'ordre de : 
1012 
dans les calculs analytiques les plus avancés. 
Dans Gabriel multiloop, pour des rapports non typiques proches de : 
1/k avec k ~= 104, 
les ancrages premiers atteignent des tailles de l'ordre de : 
P(k) ~= k5 ~= (104)5 = 1020. 
Autrement dit : 
Les premiers reconstruits par Gabriel multiloop dépassent d'un facteur d'environ 
1000 x  les hauteurs des zéros non triviaux connus de zeta(s). 
Il ne s'agit pas d'une imitation de la fonction zêta : c'est une structure parallèle, produisant
des objets arithmétiques (des nombres premiers) dans des ordres de grandeur analogues à 
ceux où la zêta produit ses zéros. 
3. Interprétation géométrique 
Dans la fonction zêta de Riemann, la structure est contrôlée par : 
s = sigma + it,  sigma = 1
2. 
La partie réelle sigma = 1/2constitue la ligne critique, où se situent les zéros non triviaux
    selon
l'hypothèse de Riemann. 
Dans Gabriel multiloop, la structure est contrôlée par une loi de puissance : 
P(k) ~= C  .  ka,  a ~= 5. 
On peut interpréter l'exposant acomme un analogue discret de la partie réelle sigmadans la 
zêta : 
-  sigma = 1/2 ->  structure des zéros, 
-  a ~= 5 ->  structure des ancrages premiers. 
Ce parallèle n'est pas une équivalence mathématique, mais une analogie structurelle : 
dans les deux cas, une exposante réelle contrôle la croissance d'objets fondamentaux 
(zéros ou premiers). 
4. Rapport typique vs rapport non typique 
Les rapports typiques (comme 1/2) produisent des ancrages faibles et réguliers. Les 
rapports non typiques (comme 1/9999) produisent des ancrages extrêmement grands. 
Cette dichotomie rappelle la distinction entre : 
-  zéros triviaux : simples, réguliers, prévisibles, 
\end{Verbatim}
\medskip
\subsubsection*{Page 48}
\begin{Verbatim}[fontsize=\footnotesize]
-  zéros non triviaux : complexes, oscillants, riches en structure. 
Dans Gabriel multiloop : 
-  les rapports typiques jouent le rôle des zéros triviaux, 
-  les rapports non typiques jouent le rôle des zéros non triviaux. 
5. Conclusion : une géométrie alternative des nombres premiers 
La structure des ancrages premiers dans Gabriel multiloop n'imite pas la fonction zêta de 
Riemann. Elle produit une géométrie polynomiale des premiers : 
P(k) ~= C  .  k5, 
qui : 
-  croît beaucoup plus rapidement que les hauteurs des zéros de zeta(s), 
-  possède une régularité remarquable, 
-  révèle une dynamique interne riche (Digamma, Zêta, offsets), 
-  et ouvre une voie entièrement nouvelle pour la compréhension des nombres 
premiers. 
 
\end{Verbatim}
\medskip

\section{Fusion des rapports typiques et non typiques dans la fonction zêta spectrale de Savard}

\subsection{Droite critique spectrale : le rapport typique $1/2$}

Dans la géométrie du spectre des nombres premiers, le rapport typique
$1/k=1/2$ (donc $k=2$) joue le rôle d'un homologue de la droite critique
de la fonction zêta de Riemann. Dans ce modèle, $n$ indique à la fois le
nombre de termes des suites $A$ et $B$ et le rang du premier associé :
\[
n=1\mapsto 2,\qquad n=2\mapsto 3,\qquad n=3\mapsto 5,\qquad
\ldots,\qquad n=49\mapsto 227.
\]

Pour les indices négatifs, on adopte la convention $P_{-m}=-p_m$, où $p_m$
est le $m$-ième premier positif :
\[
n=-1\mapsto -2,\qquad n=-2\mapsto -3,\qquad n=-3\mapsto -5,\qquad
\ldots,\qquad n=-15\mapsto -47.
\]
Ainsi, dans cette convention, les indices entiers non nuls parcourent les
premiers positifs et négatifs. Cette correspondance est une propriété du
modèle spectral présenté ici; elle n'établit pas à elle seule une équivalence
avec l'hypothèse de Riemann.

Dans la formalisation Isabelle/HOL citée dans le rapport, le théorème
\texttt{composite\_not\_prime\_i} établit qu'un entier composé ne peut être
égal à \texttt{prime\_i i}, puisque \texttt{prime\_i i} est démontré premier.
Cette propriété porte sur la fonction formalisée; le pipeline doit toujours
signaler explicitement les candidats absents ou ambigus.

\subsection{Rapports non typiques : positions et polarité spectrales}

Les rapports non typiques $1/k$, avec $k\geq 3$, conservent dans cette
interprétation la grille des positions définie par les premiers, mais $n$
représente seulement le nombre de termes des suites $A$ et $B$. Il n'est donc
pas, en général, le rang du premier retourné. La polarité associée à un
premier peut aussi dépendre du rapport considéré.

Par exemple, pour $n=10$, le rapport $1/3$ retourne $227$, qui est le
49\textsuperscript{e} premier, tandis que $1/4$ retourne $947$, qui est le
161\textsuperscript{e} premier. Selon les conventions décrites dans le
rapport, les premiers inférieurs à $881$ sont négatifs pour le rapport $1/4$.
Un même premier peut donc être positif pour un rapport non typique et négatif
pour un autre, tout en gardant son rang habituel dans la suite des premiers.
On peut interpréter les rapports non typiques comme des projections spectrales
de la structure associée au rapport typique; il s'agit ici d'une analogie du
modèle, et non d'une identité avec les zéros de la fonction zêta classique.

\subsection{Définition proposée de la fonction zêta spectrale de Savard}

L'équation proposée pour la fonction $\psi$ de Savard s'écrit
\[
\psi_{\mathrm{Savard}}(x;k,n)
=x-\frac{k^n}{S_B(k,n)}
-\log(2\pi)-\frac{1}{2}\log(1-x^{-2}),
\]
pour les valeurs où les termes de l'expression sont définis. Les cas négatifs
requièrent l'adaptation de signe décrite dans les sections précédentes.

Un ancrage premier $P$ est dit résonant dans ce modèle lorsque
$\psi_{\mathrm{Savard}}(P;k,n)$ est proche de $P$, c'est-à-dire lorsque
l'écart résiduel est faible. Les exemples présentés attribuent au rapport
typique $1/2$ une ondulation faible et décrivent une amélioration de la
proximité pour certains grands premiers. Pour les indices négatifs au-delà de
$-47$, le rapport indique que $S_B(n)$ se stabilise près de $66$; cette
observation concerne la suite négative et l'extension de l'équation à ce cas.

\subsection{Reconstruction des premiers et exclusion des composés}

La relation de reconstruction utilisée est
\[
P_i=\frac{S_B(k,n)-\mathrm{Digamma}_{\mathrm{calculé}}}{k^6}.
\]
Dans la formalisation citée, la fonction $\texttt{prime\_i}$ ne retourne que
des nombres premiers, ce qui exclut les composés de son image. L'interprétation
supplémentaire selon laquelle le remplacement d'un premier par un composé
conduirait à un indice non entier est une explication du mécanisme de
reconstruction présenté; elle doit être distinguée du théorème Isabelle
\texttt{composite\_not\_prime\_i}.

\subsection{Comptages observés et loi de croissance proposée}

Les trois catalogues dynamiques actuellement cités couvrent les plages
$k=2$ à $1001$, $1002$ à $10001$, et $10002$ à $100001$. Ils rapportent
respectivement $423$, $2604$ et $19215$ lignes marquées
\texttt{PREMIER TROUVE}, soit
\[
423+2604+19215=22242
\]
ancrages consignés sur $100000$ rapports traités. Ces nombres sont les bilans
des fichiers générés; ils ne signifient pas que chaque résultat dispose d'un
certificat HOL indépendant.

Les ancrages observés suggèrent empiriquement une croissance de type puissance :
\[
\log P(k)\approx a\log k+b,\qquad a\approx 5,
\qquad\text{soit}\qquad P(k)\approx Ck^5.
\]
Cette loi reste une approximation empirique dans les données présentées. Elle
place certains ancrages vers $10^{20}$ pour $k$ de l'ordre de $10^4$ et vers
$10^{25}$ pour $k$ de l'ordre de $10^5$.

\subsection{Interprétation zêta : analogie des exposants}

Dans la fonction zêta de Riemann, la droite critique est donnée par
$\sigma=1/2$. Dans le modèle spectral de Savard, l'exposant de croissance
$a\approx 5$ décrit la taille des ancrages premiers. On peut rapprocher ces
deux paramètres comme analogues structurels : l'un intervient dans la
description des zéros, l'autre dans la croissance des ancrages. Les deux
constructions ne sont pas mathématiquement équivalentes : la fonction zêta
classique concerne ses zéros, tandis que le modèle de Savard produit des
candidats nombres premiers.

Le rapport mentionne des hauteurs de zéros non triviaux de l'ordre de
$10^{12}$ et des ancrages pouvant atteindre l'ordre de $10^{25}$. Le quotient
des deux échelles est
\[
\frac{10^{25}}{10^{12}}=10^{13}.
\]
Cette comparaison est uniquement une comparaison d'ordres de grandeur : une
hauteur de zéro et la valeur d'un nombre premier sont des quantités de nature
différente. Faute de référence bibliographique précisant une borne vérifiable,
$10^{12}$ est présenté ici comme l'ordre de grandeur mentionné dans le rapport,
et non comme un record mondial établi.

\subsection{Hypothèse de Savard : formulation spectrale}

\textbf{Hypothèse de Savard (formulation proposée).} Pour un entier $n$ fixé
et le long des rapports non typiques $1/k$ pour lesquels le pipeline fournit
un ancrage unique, la résonance et la loi de croissance sont conjecturées sous
la forme
\[
\psi_{\mathrm{Savard}}(P_{k,n};k,n)=P_{k,n}+o(1),
\qquad
P(k)\sim Ck^5\quad (k\to\infty).
\]
Ces relations sont des hypothèses asymptotiques à vérifier; elles ne découlent
pas des seuls exemples numériques ni du théorème d'exclusion des composés.

\clearpage
\appendix
\section{Annexe --- Tables du classeur Excel}
% ======================================================================

\textbf{Classeur source :} \url{systeme_convolutif_spectral_general.xlsx}.

\subsection{Table de convolution formelle (feuille \emph{Convolution}, extrait)}

Chaque terme vaut $c_1 \cdot k^{e_1} + c_2 \cdot k^{e_2}$ : cette
représentation encode les additions et soustractions sans figer $k$.

\begin{center}\small
\begin{tabular}{@{}c c c c c c c r l@{}}
\toprule
\rowcolor{spectralgris}
Rapport & Suite & Ordre & Coeff 1 & Exp 1 & Coeff 2 & Exp 2 & Valeur & Expression formelle \\
\midrule
1/3 & A & 1 & 1 & 1 & 0 & 0 & 3 & $k^{1}$ \\
1/3 & A & 8 & 1 & 8 & 0 & 0 & 6\,561 & $k^{8}$ \\
1/3 & A & 9 & 1 & 9 & $-1$ & 7 & 17\,496 & $k^{9} - k^{7}$ \\
1/3 & A & 10 & 1 & 10 & $-1$ & 8 & 52\,488 & $k^{10} - k^{8}$ \\
1/3 & B & 9 & 1 & 9 & 0 & 0 & 19\,683 & $k^{9}$ \\
1/3 & B & 10 & 1 & 10 & $-1$ & 8 & 52\,488 & $k^{10} - k^{8}$ \\
1/3 & B & 11 & 1 & 11 & $-1$ & 9 & 157\,464 & $k^{11} - k^{9}$ \\
1/8 & A & 9 & 1 & 9 & $-1$ & 7 & 132\,120\,576 & $k^{9} - k^{7}$ \\
1/8 & A & 10 & 1 & 10 & $-1$ & 8 & 1\,056\,964\,608 & $k^{10} - k^{8}$ \\
\bottomrule
\end{tabular}
\end{center}

\subsection{Système généralisé (feuille \emph{Systeme General})}

Formules fermées exactes, valables pour tout $k$ entier $\ge 3$ et tout $n$
entier (positif ou négatif), validées contre les exemples numériques $k=3..7$ :
$\SA(k,n) = (\alpha_A(k)/2)\,k^{n} - \delta_A(k)$ et
$\SB(k,n) = (\alpha_B(k)/2)\,k^{n} - \delta_B(k)$. À $n=10$ : Digamma $= \SA
\pm A(7\ \text{ou}\ 8)$ ; les quatre possibilités servent uniquement à
déterminer l'ancrage premier du rapport. À $n \neq 10$ : rang cible $=$ rang
premier à $n=10$ $+\,n-10$ ; $P$ cible lu dans \texttt{tblPremiers} ; Digamma
calculé $= \SB(k,n) - P_{\text{cible}} \times k^{6}$. Reconstruction :
$P = [\SB(k,n) - \text{Digamma calculé}]/k^{6}$ ; un rapport sans ancrage
unique à $n=10$ demeure bloqué.

\begin{center}\scriptsize
\setlength{\tabcolsep}{2pt}
\begin{tabular}{@{}c c c r r c c r r r l@{}}
\toprule
\rowcolor{spectralgris}
Rapport & $k$ & $n$ & \multicolumn{1}{c}{Somme A} & \multicolumn{1}{c}{Somme B} & Pos. Dig. & Signe & \multicolumn{1}{c}{Digamma} & \multicolumn{1}{c}{Premier} & Rang & Statut \\
\midrule
1/3 & 3 & 10 & 79\,824 & 238\,746 & 8 & $-1$ & 73\,263 & 227 & 49 & Résolu \\
1/4 & 4 & 10 & 1\,316\,180 & 5\,260\,628 & 8 & $+1$ & 1\,381\,716 & 947 & 161 & Résolu \\
1/5 & 5 & 10 & 11\,738\,280 & 58\,675\,780 & 7 & $+1$ & 11\,816\,405 & 2\,999 & 430 & Résolu \\
1/6 & 6 & 10 & 70\,599\,858 & 423\,552\,498 & 8 & $+1$ & 72\,279\,474 & 7\,529 & 954 & Résolu \\
1/7 & 7 & 10 & 322\,966\,112 & 2\,260\,645\,142 & 8 & $-1$ & 317\,201\,311 & 16\,519 & 1\,913 & Résolu \\
1/8 & 8 & 10 & 1\,208\,259\,144 & 9\,665\,811\,016 & 8 & $-1$ & 1\,191\,481\,928 & 32\,327 & 3\,468 & Résolu \\
1/9 & 9 & 10 & 3\,874\,802\,760 & 34\,872\,693\,408 & 7 & $-1$ & 3\,870\,019\,791 & 58\,337 & 5\,906 & Résolu \\
1/10 & 10 & 10 & 11\,001\,111\,110 & 110\,010\,111\,110 & \multicolumn{5}{c}{---} & Ambigu (3 candidats) \\
1/11 & 11 & 10 & 28\,297\,321\,008 & 311\,268\,759\,538 & \multicolumn{5}{c}{---} & Ambigu (aucun candidat) \\
1/12 & 12 & 10 & 67\,080\,402\,012 & 804\,961\,838\,172 & \multicolumn{5}{c}{---} & Ambigu (2 candidats) \\
\bottomrule
\end{tabular}
\end{center}

\textbf{Cas corrigé exact --- rapport 1/8, $n = 34$ :} 32\,327 est le
3\,468\textsuperscript{ème} premier à $n=10$ ; le décalage de 24 rangs donne le
3\,492\textsuperscript{ème} premier, soit 32\,537.
$\SA(8,34) = 5\,705\,842\,489\,643\,358\,455\,620\,763\,423\,304$ ;
$\SB(8,34) = 45\,646\,739\,917\,146\,867\,644\,966\,107\,124\,296$ ;
Zêta $k^{6} = 262\,144$ ; Digamma exact
$= 45\,646\,739\,917\,146\,867\,644\,957\,577\,744\,968$ --- \emph{reconstruit
exactement}.

\subsection{Paramètres des règles et constantes des sommes fermées (feuille \emph{Parametres})}

\begin{center}\scriptsize
\setlength{\tabcolsep}{2pt}
\begin{tabular}{@{}c c c c c c l@{}}
\toprule
\rowcolor{spectralgris}
Rapport & $k$ & $n$ & Exp. reconstruction & Position Digamma & Signe Digamma & Statut \\
\midrule
1/3 & 3 & 10 & 6 & 8 & $-1$ & Ancrage $n=10$ validé ; séquence par rang active \\
1/4 & 4 & 10 & 6 & 8 & $+1$ & Ancrage $n=10$ validé ; séquence par rang active \\
1/5 & 5 & 10 & 6 & 7 & $+1$ & Ancrage $n=10$ validé ; séquence par rang active \\
1/6 & 6 & 10 & 6 & 8 & $+1$ & Ancrage $n=10$ validé ; séquence par rang active \\
1/7 & 7 & 10 & 6 & 8 & $-1$ & Ancrage $n=10$ validé ; séquence par rang active \\
1/8 & 8 & 10 & 6 & 8 & $-1$ & Ancrage $n=10$ validé ; séquence par rang active \\
1/9 & 9 & 10 & 6 & 7 & $-1$ & Ancrage $n=10$ validé ; séquence par rang active \\
1/10 & 10 & 10 & 6 & \multicolumn{2}{c}{---} & Essai-erreur requis \\
1/11 & 11 & 10 & 6 & \multicolumn{2}{c}{---} & Essai-erreur requis \\
1/12 & 12 & 10 & 6 & \multicolumn{2}{c}{---} & Essai-erreur requis \\
\bottomrule
\end{tabular}
\end{center}

Obligation universelle : chaque premier cible doit être produit une fois par le
catalogue complet. Règles actives --- $n = 10$ : A et B construites, quatre
Digamma $\pm A(7)$, $\pm A(8)$, sélection du seul entier premier, dénominateur
$k^{6}$, création de l'ancrage ; $n \neq 10$ : A et B construites, rang cible
$=$ rang base $+\,n-10$, premier lu dans \texttt{tblPremiers}, Digamma $= \Sigma B
- P \times k^{6}$, ancrage $n=10$ obligatoire.

\subsection{Moteur géo convolutif (feuille \emph{Moteur Géo Convolutif}, $k = 11$)}

Implémentation des annexes PDF 16.1--16.19 : chaque terme géométrique est le
terme entier correspondant multiplié par $g(k) = \sqrt{1+k^{2}}/k$
($g(11) = 1.0041237288352056$). Identité centrale :
$\SA^{\text{géo}} = g(k)\cdot\SA^{\text{entier}}$ ;
$\SB^{\text{géo}} = g(k)\cdot\SB^{\text{entier}}$ ;
$\text{Digamma}_{\text{géo}} = g(k)\cdot\text{Digamma}_{\text{entier}}$ ;
$\text{Zêta}_{\text{géo}} = g(k)\cdot k^{6}$ --- le facteur s'annule dans $P$.
Somme entière A $= 28\,297\,321\,008$, B $= 311\,268\,759\,538$ ; sommes
géométriques : contrôle facteur 0. Pour $k = 11$, les six branches
($A(n{-}3)\pm$, $A(n{-}2)\pm$, diagnostics $B(n{-}2)\pm$) donnent des candidats
entiers non premiers : verdict pipeline \og AUCUN P --- COMPOSÉS EXCLUS \fg{}
(6 composés exclus HOL) ; la règle standard ne sélectionne que les quatre
branches A. Référence PDF : premier imprimé 1\,611\,851 (règle A7$-$, forme
spéciale) ; correction séquentielle : ancrage rang 121\,982 à $n=10$, cible
rang 121\,987 à $n=15$ $\Rightarrow$ premier 1\,611\,901.

\subsection{Tests Digamma (feuille \emph{Tests Digamma})}

Le candidat est A ajustée $=$ Somme A $+$ signe $\cdot k^{\text{position}}$ ;
puis $P = (\text{Somme B} - \text{A ajustée})/k^{\text{exp reconstruction}}$.
Pour $1/3$ ($n=10$) : seule la combinaison position 8 / signe $-$ produit 227,
premier de rang 49 ; les trois autres candidats sont composés (221, 215, 209).
Essai convolutif $1/3$, $n = 18$ : $\SA(18) = 523\,735\,104$,
$\SB(18) = 1\,571\,204\,586$ ; les quatre options donnent des quotients
entiers, et les options position 8 produisent deux premiers consécutifs
(1\,436\,849 et 1\,436\,867) : la règle Digamma ne sélectionne pas encore une
solution unique pour $n = 18$ --- un critère supplémentaire est requis, d'où la
voie officielle du Niveau 5 (décalage de rang, Digamma calculé à rebours).

\subsection{Suites géométriques A--B (feuille \emph{Suites Géo A-B}, modèle hypoténuse)}

Pour $t = 1/k$ : typique $k=2$ ($t = 1/2$, $\SA = 2.2469862782$,
$\SB = 2.2016753304$, $\SA + \SB = 4.4486616086$) ; non typiques $k = 3..12$
(ex.\ $k=7$ : $\SA = 1.1785207100$, $\SB = 1.1784513659$). Logique de repli
automatique : Étape 1 approche algébrique ($k^{n}$ + 4 Digamma) ; si 0 premier,
Étape 2 approche géométrique (suites $\sqrt{a^{2}+b^{2}}$, extraction
floor/ceil/round de $\SA$, $\SB$, $n\SA$, $n\SB$, test de primalité).
Invariants : Digamma inchangé ; les équations pour $\forall n$ restent
identiques.

\subsection{Bilan des ancrages pour $n=10$ et extraits des grands rapports}

Les trois catalogues dynamiques couvrant successivement $k=2$--$1001$, $1002$--
$10001$, puis $10002$--$100001$ rapportent les nombres de resultats suivants.
Les nombres de la colonne "ancrages" sont les lignes marquees \texttt{PREMIER TROUVE}
par le pipeline dans les fichiers de resultats cites.

\begin{center}\small
\begin{tabular}{@{}l r r r@{}}
\toprule
\rowcolor{spectralgris}
Plage de $k$ & Rapports traites & Ancrages trouves & Proportion \\
\midrule
$2$--$1001$ & 1\,000 & 423 & 42,30\% \\
$1002$--$10001$ & 9\,000 & 2\,604 & 28,93\% \\
$10002$--$100001$ & 90\,000 & 19\,215 & 21,35\% \\
\midrule
\rowcolor{spectralazur!8}
Total ($2$--$100001$) & 100\,000 & \textbf{22\,242} & 22,242\% \\
\bottomrule
\end{tabular}
\end{center}

Les totaux reproduisent les bilans des fichiers \url{liste_k2_k1001_n10.txt},
\url{liste_k1002_k10001_n10.txt} et
\url{liste_k10002_k100001_n10.txt}. Ils decrivent les resultats consignes par
le pipeline; ils ne remplacent pas un certificat independant pour chaque ancrage.

\paragraph{Comparaison d'ordres de grandeur.}
Les ancrages des rapports proches de $1/100000$ atteignent environ $10^{25}$ dans le
catalogue ci-dessous. Le document evoque des calculs de zeros non triviaux de hauteur
de l'ordre de $10^{12}$. Le quotient de ces ordres de grandeur est
$10^{25}/10^{12}=10^{13}$, soit $10\,000$ milliards. Cette comparaison porte uniquement
sur les tailles numeriques : une hauteur de zero et la valeur d'un nombre premier sont
des quantites de nature differente. Faute de reference bibliographique precisant une
borne et un record verifiables, $10^{12}$ est presente ici comme un ordre de grandeur,
non comme un record mondial etabli.

\subsubsection*{Extrait dynamique : $k=9978$--$10001$}
\small
\begin{longtable}{@{}r >{\raggedleft\arraybackslash}p{4.3cm} c >{\raggedright\arraybackslash}p{5.2cm}@{}}
\toprule
\rowcolor{spectralgris}
Rapport & Ancre $P(n=10)$ & Etat DYN & Resultat du rapport \\
\midrule
\endfirsthead
\toprule
\rowcolor{spectralgris}
Rapport & Ancre $P(n=10)$ & Etat DYN & Resultat du rapport \\
\midrule
\endhead
\bottomrule
\endfoot
\rowcolor{spectralazur!8} $1/9978$ & $98\allowbreak\,904\allowbreak\,828\allowbreak\,370\allowbreak\,293\allowbreak\,156\allowbreak\,971$ & \textcolor{spectralbleu}{\textbf{TROUVE}} & Essais DYN : premier retourne \\
\rowcolor{spectralazur!8} $1/9979$ & $98\allowbreak\,954\allowbreak\,399\allowbreak\,755\allowbreak\,106\allowbreak\,335\allowbreak\,579$ & \textcolor{spectralbleu}{\textbf{TROUVE}} & Essais DYN : premier retourne \\
\rowcolor{spectralazur!8} $1/9980$ & $99\allowbreak\,003\allowbreak\,991\allowbreak\,013\allowbreak\,984\allowbreak\,807\allowbreak\,999$ & \textcolor{spectralbleu}{\textbf{TROUVE}} & Essais DYN : premier retourne \\
\rowcolor{spectralrouge!8} $1/9981$ & \textemdash & \textcolor{spectralrouge}{\textbf{ECHEC}} & Aucun ancrage (4 essais Digamma) \\
\rowcolor{spectralrouge!8} $1/9982$ & \textemdash & \textcolor{spectralrouge}{\textbf{ECHEC}} & Aucun ancrage (4 essais Digamma) \\
\rowcolor{spectralrouge!8} $1/9983$ & \textemdash & \textcolor{spectralrouge}{\textbf{ECHEC}} & Aucun ancrage (4 essais Digamma) \\
\rowcolor{spectralrouge!8} $1/9984$ & \textemdash & \textcolor{spectralrouge}{\textbf{ECHEC}} & Aucun ancrage (4 essais Digamma) \\
\rowcolor{spectralrouge!8} $1/9985$ & \textemdash & \textcolor{spectralrouge}{\textbf{ECHEC}} & Aucun ancrage (4 essais Digamma) \\
\rowcolor{spectralrouge!8} $1/9986$ & \textemdash & \textcolor{spectralrouge}{\textbf{ECHEC}} & Aucun ancrage (4 essais Digamma) \\
\rowcolor{spectralazur!8} $1/9987$ & $99\allowbreak\,351\allowbreak\,686\allowbreak\,808\allowbreak\,222\allowbreak\,880\allowbreak\,721$ & \textcolor{spectralbleu}{\textbf{TROUVE}} & Essais DYN : premier retourne \\
\rowcolor{spectralrouge!8} $1/9988$ & \textemdash & \textcolor{spectralrouge}{\textbf{ECHEC}} & Aucun ancrage (4 essais Digamma) \\
\rowcolor{spectralrouge!8} $1/9989$ & \textemdash & \textcolor{spectralrouge}{\textbf{ECHEC}} & Aucun ancrage (4 essais Digamma) \\
\rowcolor{spectralrouge!8} $1/9990$ & \textemdash & \textcolor{spectralrouge}{\textbf{ECHEC}} & Aucun ancrage (4 essais Digamma) \\
\rowcolor{spectralrouge!8} $1/9991$ & \textemdash & \textcolor{spectralrouge}{\textbf{ECHEC}} & Aucun ancrage (4 essais Digamma) \\
\rowcolor{spectralrouge!8} $1/9992$ & \textemdash & \textcolor{spectralrouge}{\textbf{ECHEC}} & Aucun ancrage (4 essais Digamma) \\
\rowcolor{spectralazur!8} $1/9993$ & $99\allowbreak\,650\allowbreak\,488\allowbreak\,659\allowbreak\,218\allowbreak\,583\allowbreak\,521$ & \textcolor{spectralbleu}{\textbf{TROUVE}} & Essais DYN : premier retourne \\
\rowcolor{spectralazur!8} $1/9994$ & $99\allowbreak\,700\allowbreak\,358\allowbreak\,785\allowbreak\,863\allowbreak\,712\allowbreak\,439$ & \textcolor{spectralbleu}{\textbf{TROUVE}} & Essais DYN : premier retourne \\
\rowcolor{spectralrouge!8} $1/9995$ & \textemdash & \textcolor{spectralrouge}{\textbf{ECHEC}} & Aucun ancrage (4 essais Digamma) \\
\rowcolor{spectralazur!8} $1/9996$ & $99\allowbreak\,800\allowbreak\,158\allowbreak\,937\allowbreak\,112\allowbreak\,409\allowbreak\,019$ & \textcolor{spectralbleu}{\textbf{TROUVE}} & Essais DYN : premier retourne \\
\rowcolor{spectralazur!8} $1/9997$ & $99\allowbreak\,850\allowbreak\,088\allowbreak\,973\allowbreak\,903\allowbreak\,799\allowbreak\,777$ & \textcolor{spectralbleu}{\textbf{TROUVE}} & Essais DYN : premier retourne \\
\rowcolor{spectralrouge!8} $1/9998$ & \textemdash & \textcolor{spectralrouge}{\textbf{ECHEC}} & Aucun ancrage (4 essais Digamma) \\
\rowcolor{spectralazur!8} $1/9999$ & $99\allowbreak\,950\allowbreak\,008\allowbreak\,999\allowbreak\,300\allowbreak\,019\allowbreak\,999$ & \textcolor{spectralbleu}{\textbf{TROUVE}} & Essais DYN : premier retourne \\
\rowcolor{spectralrouge!8} $1/10000$ & \textemdash & \textcolor{spectralrouge}{\textbf{ECHEC}} & Aucun ancrage (4 essais Digamma) \\
\rowcolor{spectralrouge!8} $1/10001$ & \textemdash & \textcolor{spectralrouge}{\textbf{ECHEC}} & Aucun ancrage (4 essais Digamma) \\
\end{longtable}
\normalsize

\subsubsection*{Extrait du catalogue detaille : $k=99975$--$100001$}
\small
\begin{longtable}{@{}r >{\raggedleft\arraybackslash}p{4.3cm} c >{\raggedright\arraybackslash}p{5.2cm}@{}}
\toprule
\rowcolor{spectralgris}
Rapport & Ancre $P(n=10)$ & Etat & Branche / remarque \\
\midrule
\endfirsthead
\toprule
\rowcolor{spectralgris}
Rapport & Ancre $P(n=10)$ & Etat & Branche / remarque \\
\midrule
\endhead
\bottomrule
\endfoot
\rowcolor{spectralazur!8} $1/99975$ & $9\allowbreak\,987\allowbreak\,506\allowbreak\,247\allowbreak\,438\allowbreak\,445\allowbreak\,115\allowbreak\,449\allowbreak\,949$ & \textcolor{spectralbleu}{\textbf{UNIQUE}} & Branche A(7)- \\
\rowcolor{spectralazur!8} $1/99976$ & $9\allowbreak\,988\allowbreak\,005\allowbreak\,757\allowbreak\,618\allowbreak\,485\allowbreak\,707\allowbreak\,451\allowbreak\,151$ & \textcolor{spectralbleu}{\textbf{UNIQUE}} & Branche A(7)- \\
\rowcolor{spectralgris} $1/99977$ & \textemdash & \textcolor{spectralrouge}{\textbf{ABSENT}} & -- \\
\rowcolor{spectralgris} $1/99978$ & \textemdash & \textcolor{spectralrouge}{\textbf{ABSENT}} & -- \\
\rowcolor{spectralor!20} $1/99979$ & \textemdash & \textcolor{spectralor!70!black}{\textbf{AMBIGU}} & A(8)+ : $9\,989\,504\,408\,074\,617\,108\,424\,697$; A(7)- : $9\,989\,504\,408\,074\,627\,104\,325\,117$ \\
\rowcolor{spectralgris} $1/99980$ & \textemdash & \textcolor{spectralrouge}{\textbf{ABSENT}} & -- \\
\rowcolor{spectralgris} $1/99981$ & \textemdash & \textcolor{spectralrouge}{\textbf{ABSENT}} & -- \\
\rowcolor{spectralgris} $1/99982$ & \textemdash & \textcolor{spectralrouge}{\textbf{ABSENT}} & -- \\
\rowcolor{spectralgris} $1/99983$ & \textemdash & \textcolor{spectralrouge}{\textbf{ABSENT}} & -- \\
\rowcolor{spectralgris} $1/99984$ & \textemdash & \textcolor{spectralrouge}{\textbf{ABSENT}} & -- \\
\rowcolor{spectralgris} $1/99985$ & \textemdash & \textcolor{spectralrouge}{\textbf{ABSENT}} & -- \\
\rowcolor{spectralgris} $1/99986$ & \textemdash & \textcolor{spectralrouge}{\textbf{ABSENT}} & -- \\
\rowcolor{spectralazur!8} $1/99987$ & $9\allowbreak\,993\allowbreak\,501\allowbreak\,688\allowbreak\,780\allowbreak\,714\allowbreak\,226\allowbreak\,931\allowbreak\,059$ & \textcolor{spectralbleu}{\textbf{UNIQUE}} & Branche A(8)- \\
\rowcolor{spectralgris} $1/99988$ & \textemdash & \textcolor{spectralrouge}{\textbf{ABSENT}} & -- \\
\rowcolor{spectralgris} $1/99989$ & \textemdash & \textcolor{spectralrouge}{\textbf{ABSENT}} & -- \\
\rowcolor{spectralazur!8} $1/99990$ & $9\allowbreak\,995\allowbreak\,000\allowbreak\,998\allowbreak\,900\allowbreak\,304\allowbreak\,969\allowbreak\,900\allowbreak\,999$ & \textcolor{spectralbleu}{\textbf{UNIQUE}} & Branche A(7)+ \\
\rowcolor{spectralgris} $1/99991$ & \textemdash & \textcolor{spectralrouge}{\textbf{ABSENT}} & -- \\
\rowcolor{spectralgris} $1/99992$ & \textemdash & \textcolor{spectralrouge}{\textbf{ABSENT}} & -- \\
\rowcolor{spectralgris} $1/99993$ & \textemdash & \textcolor{spectralrouge}{\textbf{ABSENT}} & -- \\
\rowcolor{spectralgris} $1/99994$ & \textemdash & \textcolor{spectralrouge}{\textbf{ABSENT}} & -- \\
\rowcolor{spectralgris} $1/99995$ & \textemdash & \textcolor{spectralrouge}{\textbf{ABSENT}} & -- \\
\rowcolor{spectralazur!8} $1/99996$ & $9\allowbreak\,998\allowbreak\,000\allowbreak\,158\allowbreak\,993\allowbreak\,720\allowbreak\,123\allowbreak\,199\allowbreak\,039$ & \textcolor{spectralbleu}{\textbf{UNIQUE}} & Branche A(7)+ \\
\rowcolor{spectralgris} $1/99997$ & \textemdash & \textcolor{spectralrouge}{\textbf{ABSENT}} & -- \\
\rowcolor{spectralgris} $1/99998$ & \textemdash & \textcolor{spectralrouge}{\textbf{ABSENT}} & -- \\
\rowcolor{spectralgris} $1/99999$ & \textemdash & \textcolor{spectralrouge}{\textbf{ABSENT}} & -- \\
\rowcolor{spectralgris} $1/100000$ & \textemdash & \textcolor{spectralrouge}{\textbf{ABSENT}} & -- \\
\rowcolor{spectralgris} $1/100001$ & \textemdash & \textcolor{spectralrouge}{\textbf{ABSENT}} & -- \\
\end{longtable}
\normalsize

Pour $k=99979$, le catalogue detaille signale deux candidats de branches differentes;
le statut reste \textbf{AMBIGU} et aucun choix unique n'est presente ici. Les Etats
\textbf{ABSENT} et \textbf{AMBIGU} ne sont donc pas comptes comme ancrages uniques dans
le bilan precedent.

\end{document}
